% Appendix C (continued): complete proofs for Exposé III, Definition III.1 to Corollary III.13.
% Sources: theory/propositions.json (statements, repaired proofs, public statements, round-2 notes),
% theory/round2_verdicts.json (suggested fixes for III.4, III.5, III.6, III.12), theory/framework.md (Exposé III),
% site/sections/06-fibres.html. The restated results follow sections/03a-fibres.tex and sections/03b-dynamics.tex.
% This file continues Appendix C and has no \section of its own.

% ====================================================================================================
\subsection{Definition III.1}\label{proof:III.1}

\begin{restate}{Definition}{III.1}{trained method; cometric; three equivalences}
A \emph{trained method} is a method together with optimizer data: a learning-rate operator $\Lambda\succ0$ (per-block
rates), an initial scale, and an update rule. Gradient flow $\dot q=-\Lambda\nabla_q(L\circ\rho)$ moves the weights by
\[\dot\theta=-K_q\nabla_\theta L,\qquad K_q=D\rho_q\,\Lambda\,D\rho_q^\top.\]
We call $K_q$ the \textbf{induced cometric}. For a backward method with compression $c_t$ and anchor $a_t$
(\thmref{Definition}{IV.1}), $K_t=a_tc_t$. Two methods are \emph{expressively equivalent} if their images agree,
\emph{first-order equivalent} if their $K_{q_0}$ agree, and \emph{dynamically equivalent} if they produce the same
$\theta$-trajectory for every loss.
\end{restate}

A definition needs no proof. This subsection fixes the conventions used in all proofs of Exposé~III and checks the two
identities that the definition records.

\paragraph{Conventions}
Weight space $\Theta=\R^{m\times n}$ carries the Frobenius inner product $\langle X,Y\rangle=\tr(X^\top Y)$, and every
parameter space carries its Euclidean (Frobenius) inner product. Gradients are identified with covectors through these
inner products, and the transpose of a linear map is its adjoint. A covector on $\Theta$ is written as a matrix, usually
$G$. When $G=\nabla_\theta L$ is a weight gradient we say so; otherwise $G$ is an arbitrary cotangent signal.

LoRA is $\rho(B,A)=W_0+sBA$ with $B\in\R^{m\times r}$, $A\in\R^{r\times n}$ and $s>0$. Its differential and the
adjoint of the differential are
\begin{equation}\label{eq:C2-lora-diff}
D\rho_{(B,A)}(\dot B,\dot A)=s(\dot BA+B\dot A),\qquad D\rho_{(B,A)}^\top(G)=(sGA^\top,\ sB^\top G),
\end{equation}
since $\langle G,s(\dot BA+B\dot A)\rangle=\langle sGA^\top,\dot B\rangle+\langle sB^\top G,\dot A\rangle$. With
$G=\nabla_\theta L(\rho(B,A))$ this gives $\nabla_B(L\circ\rho)=sGA^\top$ and $\nabla_A(L\circ\rho)=sB^\top G$, written
$\nabla_B$ and $\nabla_A$. \emph{Block-scalar rates} are $\Lambda(B,A)=(\eta_BB,\eta_AA)$ with $\eta_A,\eta_B>0$.

The gauge group $\GL_r$ acts by $g\cdot(B,A)=(Bg^{-1},gA)$, and $\rho(g\cdot q)=\rho(q)$ (\bgref{D.21}). At
$q'=g\cdot q$ the weight is unchanged, hence so is the weight gradient $G$, and \eqref{eq:C2-lora-diff} gives
\begin{equation}\label{eq:C2-gauge-grad}
\nabla_{B'}=sGA^\top g^\top=\nabla_Bg^\top,\qquad \nabla_{A'}=sg^{-\top}B^\top G=g^{-\top}\nabla_A .
\end{equation}
The \emph{full-rank locus} is $\{\rk B=\rk A=r\}$. There $BA$ has rank $r$, $U:=\operatorname{col}B=\operatorname{col}(BA)$
and $V:=\operatorname{row}A=\operatorname{row}(BA)$, and the orthogonal projectors onto them are
$P_U=B(B^\top B)^{-1}B^\top$ and $P_V=A^\top(AA^\top)^{-1}A$.

An update rule is \emph{equivariant} under $g$ (\bgref{D.18}) if, for every sequence of losses (the same data order and
dropout masks in both runs), the run started from $g\cdot q_0$ with correspondingly transformed optimizer state produces
$g\cdot q_t$ at every step $t$. Equivariance gives identical weight trajectories, since $\rho(g\cdot q_t)=\rho(q_t)$. It
is checked one step at a time. If $q_t'=g\cdot q_t$, the two runs see the same weight and the same loss, hence the same
$G$, and their gradients are related by \eqref{eq:C2-gauge-grad}. It then suffices that the steps are transported,
\begin{equation}\label{eq:C2-equivariant-step}
\delta B'=\delta B\,g^{-1},\qquad\delta A'=g\,\delta A,
\end{equation}
together with the optimizer state, and induction over $t$ concludes.

\begin{proof}[Verification of the identities in Definition III.1]
By the chain rule $\nabla_q(L\circ\rho)=D\rho_q^\top\nabla_\theta L$, so gradient flow moves the weights by
$\dot\theta=D\rho_q\dot q=-D\rho_q\Lambda D\rho_q^\top\nabla_\theta L=-K_q\nabla_\theta L$. The operator $K_q$ is
symmetric, and $\langle\gamma,K_q\gamma\rangle=\langle D\rho_q^\top\gamma,\Lambda D\rho_q^\top\gamma\rangle\ge0$. Since
$\Lambda\succ0$, $\ker K_q=\ker D\rho_q^\top$, so $\rk K_q=\rk D\rho_q\le\dim Q$. Thus $K_q$ is a cometric in the sense
of \bgref{D.23}, possibly degenerate.

A backward method updates $\theta\leftarrow\theta-\eta\,a_t(\mathrm{opt}(c_t(\nabla L)))$ (\thmref{Definition}{IV.1}).
When $\mathrm{opt}$ is the identity, that is, plain (S)GD on the compressed gradient, and $\eta$ is absorbed into $a_t$,
the continuous-time update is $\dot\theta=-a_tc_t\nabla_\theta L$, which is the statement $K_t=a_tc_t$. This operator
is symmetric positive semidefinite when $a_t=c_t^\top$, as for GaLore ($P_tP_t^\top$), LISA (a coordinate projection)
and MeZO ($z_tz_t^\top$).

Finally, for two methods trained by gradient flow, dynamical equivalence implies first-order equivalence. For a linear
loss $L(\theta)=\langle\gamma,\theta\rangle$ the initial weight velocities are $-K_{q_0}\gamma$ and $-K'_{q_0'}\gamma$,
and $\gamma$ is arbitrary.
\end{proof}

% ====================================================================================================
\subsection{Observation III.2}\label{proof:III.2}

\begin{restate}{Observation}{III.2}{chain rule for cometrics}
Let $\sigma:Q'\to Q$ and $\rho:Q\to X$ be smooth, with $X=\Theta$ or a function or output space. Train $Q'$ by gradient
flow with a constant rate operator $\Lambda'\succ0$, and put $K^\sigma_{q'}:=D\sigma_{q'}\Lambda'D\sigma_{q'}^\top$. Then
\[K^{\rho\sigma}_{q'}=D\rho_{\sigma(q')}\,K^\sigma_{q'}\,D\rho_{\sigma(q')}^\top,\]
and no metric on $Q$ is used. $K$ is a degenerate cometric (positive semidefinite, of rank $\le\dim Q'$); when $\sigma$
is not injective, $K^\sigma$ is a field along $\sigma$ and fails to be a cometric on $Q$.

For LoRA, $\rho(B,A)=W_0+sBA$ with $s=\alpha/r$ (or $\alpha/\sqrt r$ for rsLoRA), with block rates
$\Lambda=\diag(\eta_BI,\eta_AI)$ and no LoRA dropout, under gradient flow (equivalently, to first order in the step
size for GD),
\[K_{(B,A)}(G)=s^2\big(\eta_B\,GA^\top A+\eta_A\,BB^\top G\big).\]
Under Adam there is no state-independent linear cometric; \thmref{Theorem}{III.10} and \thmref{Theorem}{III.12} treat
Adam.
\end{restate}

\begin{proof}
\emph{Chain rule.} By definition $K^{\rho\sigma}_{q'}=D(\rho\sigma)_{q'}\Lambda'D(\rho\sigma)_{q'}^\top$, and
$D(\rho\sigma)_{q'}=D\rho_{\sigma(q')}D\sigma_{q'}$. Hence
\[K^{\rho\sigma}_{q'}=D\rho_{\sigma(q')}\big(D\sigma_{q'}\Lambda'D\sigma_{q'}^\top\big)D\rho_{\sigma(q')}^\top
=D\rho_{\sigma(q')}K^\sigma_{q'}D\rho_{\sigma(q')}^\top.\]
Equivalently, the flow $\dot q'=-\Lambda'\nabla_{q'}(L\circ\rho\circ\sigma)=-\Lambda'D\sigma^\top D\rho^\top dL$ moves
$x=\rho(\sigma(q'))$ by $\dot x=-D\rho\,K^\sigma D\rho^\top dL$. In this formula
$D\rho_{\sigma(q')}^\top:T^*_xX\to T^*_{\sigma(q')}Q$ and $D\sigma_{q'}^\top:T^*_{\sigma(q')}Q\to T^*_{q'}Q'$ are dual
maps, and $K^\sigma_{q'}:T^*_{\sigma(q')}Q\to T_{\sigma(q')}Q$ is built from $\Lambda'$ alone. No inner product on $Q$
enters.

\emph{Degeneracy.} As in \S\ref{proof:III.1}, $K^{\rho\sigma}_{q'}$ is symmetric and positive semidefinite, and its
rank is at most $\rk D\sigma_{q'}\le\dim Q'$. The operator $K^\sigma_{q'}$ depends on $q'$ and not only on
$\sigma(q')$. For example, with $\sigma(a,b)=ab$ from $\R^2$ to $\R$ and $\Lambda'=I$, $K^\sigma_{(a,b)}=a^2+b^2$ takes
the value $2$ at $(1,1)$ and $17/4$ at $(2,1/2)$, two points over the same value $\sigma=1$. So for non-injective
$\sigma$, $K^\sigma$ is a field along $\sigma$ (a section of $\sigma^*(TQ\otimes TQ)$) and need not descend to $Q$.

\emph{LoRA.} By \eqref{eq:C2-lora-diff}, $\Lambda D\rho^\top(G)=(\eta_BsGA^\top,\eta_AsB^\top G)$, and applying $D\rho$
gives
\[K_{(B,A)}(G)=s\big((\eta_BsGA^\top)A+B(\eta_AsB^\top G)\big)=s^2\big(\eta_BGA^\top A+\eta_ABB^\top G\big).\]
A discrete GD step $\delta B=-\eta_BsGA^\top$, $\delta A=-\eta_AsB^\top G$ changes the weight by
\[\rho(B+\delta B,A+\delta A)-\rho(B,A)=s(\delta BA+B\delta A)+s\,\delta B\,\delta A
=-K_{(B,A)}(G)+s^3\eta_A\eta_B\,GA^\top B^\top G,\]
and the last term is of order $\eta^2$. With LoRA dropout the factors are trained on masked inputs. The same computation
then gives $K_{(B,A)}$ applied to the masked signal $\tilde G$ (for one example with output gradient $\delta$, input $x$
and mask $\mathsf m$, $\tilde G=\delta(\mathsf m\odot x)^\top$), which differs from $\nabla_\theta L$. This is why the
formula with $G=\nabla_\theta L$ assumes no dropout.

\emph{Adam.} Adam's step depends on its moment buffers, so it is not a function of $q$ and $G$ alone, and even the first
step from zero state has no linear cometric. Take $B=0$ and $A$ of full row rank. Then $\nabla_A=0$, so $A$ does not
move, and the bias-corrected first step on $B$ is $\delta B=-\eta_B\phi(sGA^\top)$, with $\phi(x)=x/(\lvert x\rvert+\varepsilon)$
entrywise and $0/0=0$. (Without bias correction the step is a positive multiple of the same expression with a rescaled
$\varepsilon$.) The weight moves by $\delta W=s\,\delta B\,A=-s\eta_B\phi(sGA^\top)A$, which is nonzero whenever
$GA^\top\ne0$, because $A$ has full row rank. Suppose $\delta W=-K(G)$ for a linear map $K$. At $\varepsilon=0$,
$\phi(cx)=\phi(x)$ for $c>0$ gives $K(G)=K(cG)=cK(G)$, so $K=0$. For $\varepsilon>0$, $\lvert\phi\rvert<1$ makes
$G\mapsto\delta W$ bounded, and a bounded linear map is $0$. Both conclusions contradict $\delta W\ne0$.
\end{proof}

% ====================================================================================================
\subsection{Theorem III.3}\label{proof:III.3}

\begin{restate}{Theorem}{III.3}{gradient descent is natural exactly where cometrics agree; for LoRA's linear gauge, exactly along $\Ort(r)$}
(a) Let $\varphi:Q\to Q'$ be a diffeomorphism with $\rho'\circ\varphi=\rho$, and let $\Lambda,\Lambda'$ be the rate
operators on $Q,Q'$. Under gradient flow, the two methods induce the same $\theta$-velocity at $q$ and at $\varphi(q)$
for every loss iff $K^\rho_q=K^{\rho'}_{\varphi(q)}$, that is,
\[D\rho'_{\varphi(q)}\big(D\varphi_q\Lambda D\varphi_q^\top-\Lambda'\big)D\rho'^\top_{\varphi(q)}=0.\]
This holds whenever $D\varphi\Lambda D\varphi^\top=\Lambda'$, which is an isometry when $\Lambda=\Lambda'=I$. That
condition is necessary when $D\rho'_{\varphi(q)}$ is injective, and fails to be necessary in general. Every linear gauge
of a linear parametrisation is natural; for instance, for $\rho(x,y)=\theta_0+x+y$ the shear $(x,y)\mapsto(2x,y-x)$
preserves $\rho$ and $K$ without being orthogonal. So is every pointwise-orthogonal non-linear LoRA gauge
$\varphi(q)=\varphi_{g(q)}(q)$ with $g(q)\in\Ort(r)$.

(b) For LoRA's linear gauge $\varphi_g(B,A)=(Bg^{-1},gA)$ with block-scalar rates, at points where $A$ has full row rank
($B$ arbitrary, including $B=0$) and $r<n$, the criterion holds iff $g\in\Ort(r)$. The same holds symmetrically where $B$
has full column rank and $r<m$.
\end{restate}

\begin{proof}
(a) Fix $q\in Q$ and put $\theta=\rho(q)=\rho'(\varphi(q))$. Gradient flow with rates $\Lambda$ on $Q$ moves the weight
at $q$ by $-K^\rho_q\nabla_\theta L$ (\S\ref{proof:III.1}); gradient flow with rates $\Lambda'$ on $Q'$ moves it at
$\varphi(q)$ by $-K^{\rho'}_{\varphi(q)}\nabla_\theta L$. As $L$ ranges over the linear losses $\langle\gamma,\cdot\rangle$,
the covector $\nabla_\theta L(\theta)=\gamma$ ranges over all of $T^*_\theta\Theta$, so the two velocities agree for every
loss iff the linear maps $K^\rho_q$ and $K^{\rho'}_{\varphi(q)}$ agree. Differentiating $\rho=\rho'\circ\varphi$ gives
$D\rho_q=D\rho'_{\varphi(q)}D\varphi_q$, so
\[K^\rho_q-K^{\rho'}_{\varphi(q)}=D\rho'_{\varphi(q)}\big(D\varphi_q\Lambda D\varphi_q^\top-\Lambda'\big)D\rho'^\top_{\varphi(q)},\]
which is the displayed criterion.

If $D\varphi_q\Lambda D\varphi_q^\top=\Lambda'$, the bracket vanishes. For $\Lambda=\Lambda'=I$ the condition says that
$D\varphi_q$ is orthogonal; in general it says that $\varphi$ is an isometry from $(Q,\langle\cdot,\Lambda^{-1}\cdot\rangle)$
to $(Q',\langle\cdot,\Lambda'^{-1}\cdot\rangle)$. If $D\rho'_{\varphi(q)}$ is injective, it has a left inverse $E$, and
multiplying $D\rho'MD\rho'^\top=0$ by $E$ on the left and by $E^\top$ on the right gives $M=0$; so the condition is then
necessary.

It is not necessary in general. Let $\rho(x)=\theta_0+\mathcal Lx$ be affine with linear part $\mathcal L$, and let $P$
be a linear gauge, $\mathcal LP=\mathcal L$, with $\Lambda'=\Lambda$. Then $D\rho'=\mathcal L$ and $D\varphi=P$, and
$\mathcal LP\Lambda P^\top\mathcal L^\top=\mathcal L\Lambda\mathcal L^\top$, so the criterion holds for every
$\Lambda$. For $\rho(x,y)=\theta_0+x+y$ and $\Lambda=I$, the shear $P(x,y)=(2x,y-x)$ satisfies $\mathcal LP=\mathcal L$,
while in block form $PP^\top=\begin{psmallmatrix}4&-2\\-2&2\end{psmallmatrix}\ne I$.

For a pointwise gauge $\varphi(q)=\varphi_{g(q)}(q)$ of LoRA, with $g:Q\to\GL_r$ smooth and $\varphi$ a diffeomorphism,
the product rule gives $D\varphi_q=\varphi_{g(q)}+V_q$. Here $\varphi_{g(q)}$ is the linear map
$(\dot B,\dot A)\mapsto(\dot Bg(q)^{-1},g(q)\dot A)$, and $V_q(v)=\frac{d}{d\epsilon}\big|_{\epsilon=0}\varphi_{g(q+\epsilon v)}(q)$
is tangent to the $\GL_r$-orbit of $\varphi(q)$. That orbit lies in a fibre of $\rho$, so $D\rho_{\varphi(q)}V_q=0$ and
$D\rho_{\varphi(q)}D\varphi_q=D\rho_{\varphi(q)}\varphi_{g(q)}$. The criterion depends on $D\varphi_q$ only through
$D\rho_{\varphi(q)}D\varphi_q$, so it reduces to the criterion for the linear gauge $\varphi_{g(q)}$ at $q$, which holds
for $g(q)\in\Ort(r)$ by (b).

(b) Here $Q'=Q$, $\rho'=\rho$, $\Lambda'=\Lambda$ is block-scalar and $\varphi_g(q)=g\cdot q$. By
\thmref{Observation}{III.2},
\[K_{g\cdot q}(G)=s^2\big(\eta_BGA^\top g^\top gA+\eta_ABg^{-1}g^{-\top}B^\top G\big),\]
so $K_q=K_{g\cdot q}$ iff, for every $G\in\R^{m\times n}$,
\[GX=YG,\qquad X=\eta_BA^\top(g^\top g-I)A\in\R^{n\times n},\quad Y=\eta_AB(I-g^{-1}g^{-\top})B^\top\in\R^{m\times m}.\]

\emph{Matrix-unit lemma: if $GX=YG$ for every $G\in\R^{m\times n}$, then $X=cI_n$ and $Y=cI_m$ for some $c\in\R$.} Take
$G=E_{ij}=e_ie_j^\top$. Comparing the $(k,l)$ entries of $E_{ij}X$ and $YE_{ij}$ gives $\delta_{ik}X_{jl}=Y_{ki}\delta_{jl}$.
With $k=i$ and $l\ne j$ this gives $X_{jl}=0$; with $l=j$ and $k\ne i$ it gives $Y_{ki}=0$; with $k=i$ and $l=j$ it gives
$X_{jj}=Y_{ii}$ for all $i,j$.

Now $\rk X\le r<n$, so $X=cI_n$ forces $c=0$, that is, $A^\top(g^\top g-I)A=0$. Since $A$ has full row rank, $AA^\top$ is
invertible, and multiplying by $(AA^\top)^{-1}A$ on the left and by $A^\top(AA^\top)^{-1}$ on the right gives
$g^\top g=I$. Conversely, $g\in\Ort(r)$ makes $X$ and $Y$ vanish. Only the full row rank of $A$ and $r<n$ were used, and
$B$ is arbitrary, including $B=0$. If instead $B$ has full column rank and $r<m$, then $\rk Y\le r<m$ forces $c=0$, so
$B(I-g^{-1}g^{-\top})B^\top=0$; multiplying by $(B^\top B)^{-1}B^\top$ on the left and by $B(B^\top B)^{-1}$ on the right
gives $g^{-1}g^{-\top}=I$, and again $g\in\Ort(r)$.
\end{proof}

\begin{remark*}{the family of plain-GD dynamics; LoRA+ under GD}
The formula for $K_{g\cdot q}$ depends on $g$ only through $g^\top g$, since $g^{-1}g^{-\top}=(g^\top g)^{-1}$, and
$g^\top g=h^\top h$ iff $hg^{-1}\in\Ort(r)$. So along a gauge orbit the cometrics of plain gradient descent are indexed by
$\Ort(r)\backslash\GL_r$, which $g\mapsto g^\top g$ identifies with $\Sym_r^+$ (\bgref{D.27}, \auxref{B.19}(a)). For $g=cI$ with $c>0$,
gradient flow with a single rate $\eta$ started from $g\cdot q$, pulled back by $g^{-1}$, satisfies $\dot B=c\dot B'=-\eta c^2\nabla_B$
and $\dot A=c^{-1}\dot A'=-\eta c^{-2}\nabla_A$ by \eqref{eq:C2-gauge-grad}. This is LoRA with rates
$(\eta_A,\eta_B)=(\eta c^{-2},\eta c^2)$, and the same holds step by step for GD, whose steps are linear in the gradients.
Choosing $c=\lambda^{1/4}$ and $\eta=\sqrt\lambda\,\eta_A$ gives LoRA+ with ratio $\lambda=\eta_B/\eta_A$. Under gradient
descent, LoRA+ is therefore the scalar gauge $\lambda^{1/4}I$ together with an overall change of rate.
\end{remark*}

% ====================================================================================================
\subsection{Proposition III.4}\label{proof:III.4}

\begin{restate}{Proposition}{III.4}{cometrics that descend to the rank-$r$ manifold; known in part}
Write $U=\operatorname{col}B$ and $V=\operatorname{row}A$. Statements (a), (b) and the first sentence of (d) are on the
full-rank locus; (c) needs only $\rk A_0=r$, with $B$ arbitrary, including $B_0=0$. All statements concern gradient
flow, or (S)GD to first order, on the stated gradients, undamped and without Adam. Cometrics include the learning rate.

(a) Scaled GD, $\delta B=-\eta\nabla_B(AA^\top)^{-1}$ and $\delta A=-\eta(B^\top B)^{-1}\nabla_A$, is $\GL_r$-equivariant
with first-order cometric $K(G)=\eta s^2(P_UG+GP_V)$.

(b) LoRA-Pro's equivalent gradient $\tilde g=P_UG+(I-P_U)GP_V$ is the Frobenius-orthogonal projection
$\Pi_{T_{W-W_0}\MM_r}(G)$. The flow of its SGD variant (LoRA-Pro Alg.~1) is the Riemannian gradient flow on
$W_0+\MM_r$.

(c) LoRA-FA's corrected gradient $s^{-2}\nabla_B(A_0A_0^\top)^{-1}$, fed to SGD with $A_0$ frozen, gives
$K(G)=\eta\,GP_V$ with $V=\operatorname{row}A_0$, which depends on $A_0$ only through $V\in\Gr(r,n)$. Its original form,
plain SGD on $B$, has $K(G)=s^2\eta_BGA_0^\top A_0$.

(d) The cometrics of (a) and (b) are functions of $W$ on the full-rank locus and define vector fields tangent to
$W_0+\MM_r$. The field of (c) is $G\mapsto\eta GP_V$ with $V=\operatorname{row}A_0$. Its flow is the Frobenius gradient
flow on the slice $W_0+\{XA_0\}=W_0+\R^{m\times n}P_V$, of dimension $mr$, which contains $W_0$ and lies in $W_0+\Mr$.
The rank-$r$ part of the slice, $N_V=\{\operatorname{row}(W-W_0)=V\}$, is open and dense in it and is a submanifold of
$W_0+\MM_r$. On $N_V$ the field is tangent to $\MM_r$ and equals $\eta GP_{\operatorname{row}(W-W_0)}$, a function of
$W$. So (c) descends too, as a degenerate cometric of rank $mr$ (less than $r(m+n-r)$ when $r<n$) whose flow keeps
$\operatorname{row}(W-W_0)=\operatorname{row}A_0$ fixed. The cometric of plain GD depends on the representative; it is
invariant only under $\Ort(r)$.
\end{restate}

\begin{proof}
In (a), (b) and the first sentence of (d), $(B,A)$ lies on the full-rank locus and the conventions of
\S\ref{proof:III.1} apply. Write $\mathcal T:=T_{W-W_0}\MM_r$, where $W-W_0=sBA$ has rank $r$, column space $U$ and row
space $V$.

(a) The steps are $\delta B=-\eta\nabla_B(AA^\top)^{-1}=-\eta s\,GA^\top(AA^\top)^{-1}$ and
$\delta A=-\eta(B^\top B)^{-1}\nabla_A=-\eta s\,(B^\top B)^{-1}B^\top G$. Hence
\begin{align*}
\delta W&=s(\delta BA+B\delta A)+s\,\delta B\,\delta A=-\eta s^2\big(GA^\top(AA^\top)^{-1}A+B(B^\top B)^{-1}B^\top G\big)+O(\eta^2)\\
&=-\eta s^2(GP_V+P_UG)+O(\eta^2),
\end{align*}
and the flow version gives $\dot W=-\eta s^2(P_UG+GP_V)$ exactly. Left multiplication by $P_U$ and right
multiplication by $P_V$ are orthogonal projections of $(\R^{m\times n},\langle\cdot,\cdot\rangle)$, so
$K(G)=\eta s^2(P_UG+GP_V)$ is symmetric positive semidefinite. For equivariance, let $q'=g\cdot q$. By
\eqref{eq:C2-gauge-grad} and $(A'A'^\top)^{-1}=(gAA^\top g^\top)^{-1}=g^{-\top}(AA^\top)^{-1}g^{-1}$,
\[\delta B'=-\eta\nabla_Bg^\top g^{-\top}(AA^\top)^{-1}g^{-1}=\delta B\,g^{-1},\]
and $(B'^\top B')^{-1}=(g^{-\top}B^\top Bg^{-1})^{-1}=g(B^\top B)^{-1}g^\top$ gives
$\delta A'=-\eta g(B^\top B)^{-1}g^\top g^{-\top}\nabla_A=g\,\delta A$. So the discrete step satisfies
\eqref{eq:C2-equivariant-step} for every $g\in\GL_r$.

(b) \emph{The tangent space.} We show that
\[\{XA+BY:\ X\in\R^{m\times r},\ Y\in\R^{r\times n}\}=\{D:\ P_{U^\perp}DP_{V^\perp}=0\}=\mathcal T.\]
The first inclusion $\subseteq$ holds because $AP_{V^\perp}=0$ and $P_{U^\perp}B=0$. Conversely, if
$P_{U^\perp}DP_{V^\perp}=0$, then
$D=P_UD+P_{U^\perp}DP_V=B\big[(B^\top B)^{-1}B^\top D\big]+\big[P_{U^\perp}DA^\top(AA^\top)^{-1}\big]A$. The common space
has dimension $mn-(m-r)(n-r)=r(m+n-r)=\dim\MM_r$ (\bgref{C.12}). The curves $t\mapsto s(B+t\dot B)(A+t\dot A)$ lie in
$\MM_r$ for small $t$, so $\{\dot BA+B\dot A\}\subseteq\mathcal T$, and equality of dimensions gives equality.

\emph{The projection.} Let $\mathcal N=\{P_{U^\perp}EP_{V^\perp}:E\in\R^{m\times n}\}$. For $D\in\mathcal T$ we
have $\langle D,P_{U^\perp}EP_{V^\perp}\rangle=\langle P_{U^\perp}DP_{V^\perp},E\rangle=0$, so
$\mathcal N\perp\mathcal T$. Moreover $G=(G-P_{U^\perp}GP_{V^\perp})+P_{U^\perp}GP_{V^\perp}$, with the first term in
$\mathcal T$ and the second in $\mathcal N$. Hence
\[\Pi_{\mathcal T}(G)=G-(I-P_U)G(I-P_V)=P_UG+GP_V-P_UGP_V=P_UG+(I-P_U)GP_V=\tilde g.\]

\emph{LoRA-Pro's SGD variant.} Algorithm~1 of \citet{wang2024lorac} takes the steps $\delta B=-\eta g^B$ and
$\delta A=-\eta g^A$ along a pair that minimises $\lVert s(g^BA+Bg^A)-G\rVert_F$ (their Theorem~2.1). As $(g^B,g^A)$
ranges over all pairs, $s(g^BA+Bg^A)$ ranges over $\mathcal T$, and the closest point of $\mathcal T$ to $G$ is
$\Pi_{\mathcal T}G$. So every minimiser satisfies $s(g^BA+Bg^A)=\Pi_{\mathcal T}G$; explicitly, the minimisers are
\[g^A=s^{-1}(B^\top B)^{-1}B^\top G+XA,\qquad g^B=s^{-1}(I-P_U)GA^\top(AA^\top)^{-1}-BX,\qquad X\in\R^{r\times r},\]
for which $sBg^A=P_UG+sBXA$ and $sg^BA=(I-P_U)GP_V-sBXA$. For every choice of the free matrix $X$,
\[\delta W=s(\delta BA+B\delta A)+s\,\delta B\,\delta A=-\eta\,\Pi_{\mathcal T}G+O(\eta^2).\]
The flow version is $\dot W=-\eta\,\Pi_{T_{W-W_0}\MM_r}\nabla_\theta L(W)$. Since $\Pi_{\mathcal T}\nabla_\theta L$ is the
Riemannian gradient of $L$ restricted to $W_0+\MM_r$ for the metric induced by the Frobenius inner product
(\bgref{D.24}, \auxref{B.30}; \citealp[Ch.~3]{xabsil2008optimization}), this is the Riemannian gradient flow on $W_0+\MM_r$, run at
rate $\eta$.

(c) Here $A=A_0$ is frozen with $\rk A_0=r$, and $B$ is arbitrary. The corrected step is
$\delta B=-\eta s^{-2}\nabla_B(A_0A_0^\top)^{-1}=-\eta s^{-1}GA_0^\top(A_0A_0^\top)^{-1}$, so
\[\Delta W=s\,\delta B\,A_0=-\eta\,GA_0^\top(A_0A_0^\top)^{-1}A_0=-\eta\,GP_V,\]
exactly, since $A_0$ does not move. $P_V$ depends only on $V=\operatorname{row}A_0$. The original LoRA-FA step
$\delta B=-\eta_B\nabla_B=-\eta_BsGA_0^\top$ gives $\Delta W=s\,\delta B\,A_0=-s^2\eta_BGA_0^\top A_0$.

(d) On the full-rank locus $P_U$ and $P_V$ are the projectors onto $\operatorname{col}(W-W_0)$ and
$\operatorname{row}(W-W_0)$, so the cometrics of (a) and (b) are functions of $W$. Their values lie in $\mathcal T$:
this is clear for $\Pi_{\mathcal T}$, and $P_{U^\perp}(P_UG+GP_V)P_{V^\perp}=0$.

For (c), the trajectory satisfies $W-W_0=sBA_0\in\{XA_0:X\in\R^{m\times r}\}$, and $\{XA_0\}=\{Z:\operatorname{row}Z\subseteq V\}=\R^{m\times n}P_V$,
a linear space of dimension $mr$. Right multiplication by $P_V$ is the Frobenius-orthogonal projection onto it, so
$\dot W=-\eta\nabla_\theta L(W)P_V$ is the gradient flow of $L$ restricted to the affine slice $W_0+\R^{m\times n}P_V$
with its flat Frobenius metric. The slice contains $W_0$ ($X=0$) and lies in $W_0+\Mr$, since $\rk XA_0\le r$. Because
$X\mapsto XA_0$ is injective, $\rk XA_0=\rk X$, and $\operatorname{row}(XA_0)=V$ iff $\rk X=r$. So
$N_V=W_0+\{XA_0:\rk X=r\}$, which is open in the slice because rank $r$ is an open condition (\auxref{B.1}(a)), and
dense because every $X\in\R^{m\times r}$ ($m\ge r$) is a limit of rank-$r$ matrices $X+tE$ as $t\to0$: if $E$ has
$I_r$ in its first $r$ rows and zeros elsewhere, the top $r\times r$ minor of $X+tE$ is a monic polynomial of degree $r$
in $t$, so it vanishes for only finitely many $t$. $N_V$ is an open
subset of an affine subspace, hence an embedded submanifold of $\R^{m\times n}$, and it lies in the embedded
submanifold $W_0+\MM_r$. The inclusion $N_V\hookrightarrow W_0+\MM_r$ is smooth by \auxref{B.2}(c), an immersion
because its composite with the inclusion $W_0+\MM_r\hookrightarrow\R^{m\times n}$ is one, and a homeomorphism onto its
image because both sets carry the subspace topology of $\R^{m\times n}$. So it is an embedding, and $N_V$ is an embedded
submanifold of $W_0+\MM_r$ (\bgref{C.4}). At $W\in N_V$, $\operatorname{row}(W-W_0)=V$, so $GP_V=GP_{\operatorname{row}(W-W_0)}$ is a
function of $W$, and, with $U_W=\operatorname{col}(W-W_0)$, $P_{U_W^\perp}(GP_V)P_{V^\perp}=0$ shows that it lies in $T_{W-W_0}\MM_r$ (it lies in
$\R^{m\times n}P_V=T N_V$ as well). The operator $G\mapsto\eta GP_V$ has rank $\dim\R^{m\times n}P_V=mr$, and
$mr<r(m+n-r)$ when $r<n$. Along the flow $W(t)$ stays in the slice, so $\operatorname{row}(W(t)-W_0)\subseteq V$ for
all $t$, with equality whenever $W(t)\in N_V$.

The last sentence of (d) is \thmref{Theorem}{III.3}(b).
\end{proof}

\begin{remark*}{the departures listed under Practice}
The symmetry statements in the Practice paragraph after \thmref{Proposition}{III.4} are proved with
\thmref{Theorem}{III.12}: damping $(B^\top B+\delta I)^{-1}$, $(AA^\top+\delta I)^{-1}$ with $\delta>0$ leaves exactly
$\Ort(r)$; an AdamW step applied after the preconditioner, as in scaled AdamW and in HF's LoRA-FA optimizer, leaves
exactly the signed permutations $B_r$; and LoRA-Pro's Sylvester-optimal free matrix leaves exactly $\Ort(r)$ at the
parameter level. For LoRA-FA with the undamped correction and AdamW at $\varepsilon=0$, the first step from $B_0=0$ is
$\Delta W=-s\eta_B\operatorname{sign}\big(GA_0^\top(A_0A_0^\top)^{-1}\big)A_0$. Replacing $A_0$ by $2A_0$, which has the
same row space, doubles this step, so the trajectory depends on $A_0$ beyond its row space (see also
\S\ref{proof:III.13}).
\end{remark*}

% ====================================================================================================
\subsection{Theorem III.5}\label{proof:III.5}

\begin{restate}{Theorem}{III.5}{Noether charges and isometries for cotangent signals; known, \citet{zhao2022symmetries}}
Let a Lie group act linearly on $Q=\R^N$ with $\rho\circ k=\rho$, let $\Lambda\succ0$ be constant, and let
$\gamma(t)\in T^*_{\rho(q(t))}\Theta$ be any cotangent signal, gradient or not. Consider the continuous-time updates
\[\dot q=-\Lambda\,D\rho_q^\top\gamma(t).\]

(i) For every generator $\xi$ with $\Lambda^{-1}\xi$ symmetric, the charge $J_\xi(q)=\tfrac12\langle q,\Lambda^{-1}\xi q\rangle$
is conserved.

(ii) For every generator $\xi$ with $\Lambda^{-1}\xi$ antisymmetric, $k=e^{\tau\xi}$ preserves
$\langle\cdot,\Lambda^{-1}\cdot\rangle$, so $K_{kq}=K_q$ along the orbit; for the same cotangent $\gamma$, the
$\theta$-velocities at $q$ and at $kq$ agree. Suppose moreover that $\gamma$ depends on $q$ only through $(t,\rho(q))$, as
$\gamma=\nabla_\theta L$ or an open-loop signal does, and that solutions are unique ($\gamma$ measurable in $t$ and
locally Lipschitz in $\theta$, $\rho$ of class $C^{1,1}_{\mathrm{loc}}$). Then the $\theta$-trajectories from $q_0$ and
from $kq_0$ coincide. If $\gamma$ uses Euclidean norms on $Q$, as global-norm clipping does, the trajectory statement
holds when, in addition, $k$ is a Euclidean isometry, and can fail otherwise. So (i) and the velocity statement of (ii)
hold for arbitrary $\gamma$, while the trajectory statement needs a gradient-like signal.

(iii) When $\mathfrak g$ is closed under $\xi\mapsto\Lambda\xi^\top\Lambda^{-1}$, every generator splits into a part of
type (i) and a part of type (ii), a Cartan decomposition $\mathfrak g=\mathfrak k\oplus\mathfrak p$. This holds for
LoRA's $\gl_r$ with per-block scalar rates. Without closure, there are generators of neither type. With per-rank-channel
rates $\Lambda(B,A)=(BD_B,D_AA)$, where $D_B=\diag(\eta_{B,j})$ and $D_A=\diag(\eta_{A,j})$, the generator $\xi_X$ is of
type (i) iff $XD_B^{-1}$ and $D_A^{-1}X$ are both symmetric, and of type (ii) iff both are antisymmetric. An off-diagonal
pair $(i,j)$ contributes one charge and one isometry iff $\eta_{B,i}\eta_{A,i}=\eta_{B,j}\eta_{A,j}$, and $\mathfrak g$
is closed iff all the products $\eta_{B,j}\eta_{A,j}$ are equal. In that case one gets a $D$-conjugate of
$\so_r\oplus\Sym_r$ whose isometries are not all Euclidean unless the rates are
block-scalar, so the clipping caveat of (ii) applies. For pairwise distinct
products, only the diagonal $X$ give charges and no nonzero $X$ gives an isometry.
\end{restate}

\begin{proof}
Part (i) is Noether's theorem in the form of \auxref{B.29}, which allows an arbitrary cotangent signal; we give the
short argument in full. Solutions are absolutely continuous and the identities below hold for almost
every $t$. For $\xi$ in the Lie algebra
$\mathfrak g\subseteq\gl(N)$ of the group, $\rho(e^{\tau\xi}q)=\rho(q)$ for all $\tau$; differentiating at $\tau=0$
gives the infinitesimal invariance
\begin{equation}\label{eq:C2-infinitesimal}
D\rho_q(\xi q)=0\qquad(q\in Q,\ \xi\in\mathfrak g).
\end{equation}

(i) If $\Lambda^{-1}\xi$ is symmetric, the gradient of $J_\xi$ is $\Lambda^{-1}\xi q$. Along a solution, using that
$\Lambda$ is symmetric and \eqref{eq:C2-infinitesimal},
\[\tfrac{d}{dt}J_\xi(q)=\langle\Lambda^{-1}\xi q,\dot q\rangle=-\langle\Lambda^{-1}\xi q,\Lambda D\rho_q^\top\gamma\rangle
=-\langle\xi q,D\rho_q^\top\gamma\rangle=-\langle D\rho_q\xi q,\gamma\rangle=0.\]
Nothing about $\gamma$ was used, so it may depend on $q$ in any way.

(ii) \emph{Isometry and cometric.} If $\Lambda^{-1}\xi$ is antisymmetric, then $\xi^\top\Lambda^{-1}=-\Lambda^{-1}\xi$, so
\[\tfrac{d}{d\tau}\big(e^{\tau\xi^\top}\Lambda^{-1}e^{\tau\xi}\big)=e^{\tau\xi^\top}(\xi^\top\Lambda^{-1}+\Lambda^{-1}\xi)e^{\tau\xi}=0,\]
and $k=e^{\tau\xi}$ satisfies $k^\top\Lambda^{-1}k=\Lambda^{-1}$, equivalently $k\Lambda k^\top=\Lambda$. From
$\rho\circ k=\rho$ we get $D\rho_q=D\rho_{kq}k$, so
\[K_q=D\rho_{kq}\,k\Lambda k^\top D\rho_{kq}^\top=D\rho_{kq}\Lambda D\rho_{kq}^\top=K_{kq}.\]
Since $\rho(kq)=\rho(q)$, the same covector $\gamma$ acts at both points, and the velocities $-K_q\gamma$ and
$-K_{kq}\gamma$ agree. The rest of the argument uses only the two properties $k\Lambda k^\top=\Lambda$ and $\rho\circ k=\rho$,
so it applies to every such $k$, not only to those of the form $e^{\tau\xi}$.

\emph{Trajectories.} Let $\gamma=\gamma(t,\rho(q))$ and let $q(t)$ solve $\dot q=F(t,q):=-\Lambda D\rho_q^\top\gamma(t,\rho(q))$
from $q_0$. Put $p(t)=kq(t)$. Then, using $D\rho_q^\top=k^\top D\rho_{kq}^\top$ and $\rho(p)=\rho(q)$,
\[\dot p=k\dot q=-k\Lambda k^\top D\rho_p^\top\gamma(t,\rho(q))=-\Lambda D\rho_p^\top\gamma(t,\rho(p))=F(t,p),\]
so $p$ solves the same equation from $kq_0$. Under the stated regularity, read locally uniformly in $t$, $F$ is measurable in
$t$ and locally bounded and locally Lipschitz in $p$, because $D\rho$ and $\rho$ are locally Lipschitz and $\gamma$ is
locally Lipschitz in $\theta$. By Gronwall's inequality two solutions with the same initial value coincide (\auxref{B.28}), so $p$ is the solution from
$kq_0$, and $\rho(p(t))=\rho(q(t))$ for all $t$.

\emph{Clipping.} Suppose $\gamma$ depends on $q$ through $(t,\rho(q))$ and through Euclidean norms
$\lVert D\rho_q^\top v\rVert$ of pulled-back covectors $v$ that are themselves functions of $(t,\rho(q))$. Global-norm
clipping is of this kind:
\[\gamma=\min\big(1,C/\lVert D\rho_q^\top\nabla_\theta L\rVert\big)\,\nabla_\theta L .\]
If $k$ is also a Euclidean isometry, then
$\lVert D\rho_{kq}^\top v\rVert=\lVert k^{-\top}D\rho_q^\top v\rVert=\lVert D\rho_q^\top v\rVert$, so $\gamma$ takes the
same value at $q$ and at $kq$, and the argument above goes through. Otherwise it can fail. Take $Q=\R^2$,
$\Theta=\R$, $\Lambda=\diag(1,4)$, $\rho(q)=\tfrac12q^\top\Lambda^{-1}q$ and $L(\theta)=\theta^2/2$. The group
$\{k:k^\top\Lambda^{-1}k=\Lambda^{-1}\}$ acts linearly and preserves $\rho$; its generator $\xi=\Lambda^{1/2}J\Lambda^{-1/2}$,
with $J=\begin{psmallmatrix}0&-1\\1&0\end{psmallmatrix}$, has $\Lambda^{-1}\xi=\Lambda^{-1/2}J\Lambda^{-1/2}$
antisymmetric, and $k=e^{(\pi/2)\xi}=\Lambda^{1/2}e^{(\pi/2)J}\Lambda^{-1/2}$ maps $q_0=(1,0)$ to $kq_0=(0,2)$, both with
$\theta=1/2$. Here $D\rho_q^\top\gamma=\gamma\Lambda^{-1}q$ and $\dot\theta=-\gamma\,q^\top\Lambda^{-1}q=-2\theta\gamma$.
The parameter gradients $L'(\theta)\Lambda^{-1}q$ have norms $1/2$ at $q_0$ and $1/4$ at $kq_0$, so with $C=1/10$ the
clipping factors are $1/5$ and $2/5$, and the initial velocities are $\dot\theta=-1/10$ and $\dot\theta=-1/5$. The
$\theta$-trajectories differ, while without clipping they coincide by the argument above.

(iii) \emph{Splitting.} Put $\vartheta(\xi)=-\Lambda\xi^\top\Lambda^{-1}$. It is an involution of $\gl(N)$
($\vartheta^2(\xi)=\Lambda\Lambda^{-1}\xi\Lambda\Lambda^{-1}=\xi$) and a Lie algebra automorphism:
$\vartheta([\xi,\zeta])=-\Lambda(\zeta^\top\xi^\top-\xi^\top\zeta^\top)\Lambda^{-1}=[\vartheta\xi,\vartheta\zeta]$. A
generator is of type (ii) iff $\Lambda^{-1}\xi$ is antisymmetric, iff $\Lambda\xi^\top\Lambda^{-1}=-\xi$, iff
$\vartheta\xi=\xi$; it is of type (i) iff $\vartheta\xi=-\xi$. If $\mathfrak g$ is closed under
$\xi\mapsto\Lambda\xi^\top\Lambda^{-1}$, it is $\vartheta$-stable, and
\[\xi=\tfrac12(\xi+\Lambda\xi^\top\Lambda^{-1})+\tfrac12(\xi-\Lambda\xi^\top\Lambda^{-1})\]
writes every $\xi\in\mathfrak g$ as a type (i) part (in the $-1$ eigenspace $\mathfrak p$ of $\vartheta$) plus a type (ii)
part (in the $+1$ eigenspace $\mathfrak k$). Since $\vartheta$ is an automorphism, $[\mathfrak k,\mathfrak k]\subseteq\mathfrak k$,
$[\mathfrak k,\mathfrak p]\subseteq\mathfrak p$ and $[\mathfrak p,\mathfrak p]\subseteq\mathfrak k$, which is the Cartan
decomposition (\bgref{D.27}, \auxref{B.19}(b)). If some $\xi\in\mathfrak g$ has $\Lambda\xi^\top\Lambda^{-1}\notin\mathfrak g$, then
$\vartheta\xi\ne\pm\xi$, so $\xi$ is of neither type.

\emph{LoRA.} The generators of the gauge are $\xi_X(B,A)=(BX,-XA)$ for $X\in\gl_r$ (the derivative of $\tau\mapsto e^{-\tau X}\cdot(B,A)$),
and $X\mapsto\xi_X$ is injective. For per-rank-channel rates $\Lambda(B,A)=(BD_B,D_AA)$,
\[\Lambda^{-1}\xi_X(B,A)=(BXD_B^{-1},\ -D_A^{-1}XA),\qquad (\Lambda^{-1}\xi_X)^\top(B,A)=\big(B(XD_B^{-1})^\top,\ -(D_A^{-1}X)^\top A\big),\]
because $\langle B',BM\rangle=\langle B'M^\top,B\rangle$ and $\langle A',NA\rangle=\langle N^\top A',A\rangle$. Hence
$\Lambda^{-1}\xi_X$ is symmetric iff $XD_B^{-1}$ and $D_A^{-1}X$ are symmetric, and antisymmetric iff both are
antisymmetric. Entrywise, $(XD_B^{-1})_{ij}=X_{ij}/\eta_{B,j}$ and $(D_A^{-1}X)_{ij}=X_{ij}/\eta_{A,i}$. For $i\ne j$, the
two symmetry conditions read $X_{ji}=X_{ij}\eta_{B,i}/\eta_{B,j}$ and $X_{ji}=X_{ij}\eta_{A,j}/\eta_{A,i}$; the
antisymmetry conditions are the same with a minus sign. Both admit $X_{ij}\ne0$ iff
$\eta_{B,i}/\eta_{B,j}=\eta_{A,j}/\eta_{A,i}$, that is, iff $\eta_{B,i}\eta_{A,i}=\eta_{B,j}\eta_{A,j}$, and then the pair
$(i,j)$ carries a one-dimensional space of each type. On the diagonal the symmetry conditions are empty and the
antisymmetry conditions force $X_{ii}=0$. So for pairwise distinct products the type (i) generators are exactly the
diagonal $X$, and the only type (ii) generator is $X=0$.

For closure, $\xi_X^\top(B,A)=(BX^\top,-X^\top A)$, so $\Lambda\xi_X^\top\Lambda^{-1}(B,A)=(BD_B^{-1}X^\top D_B,\ -D_AX^\top D_A^{-1}A)$.
An operator $(B,A)\mapsto(BM,-NA)$ lies in $\mathfrak g=\{\xi_Y\}$ iff $M=N$, so closure holds iff
$D_B^{-1}X^\top D_B=D_AX^\top D_A^{-1}$ for all $X$. The $(i,j)$ entries are $X_{ji}\eta_{B,j}/\eta_{B,i}$ and
$X_{ji}\eta_{A,i}/\eta_{A,j}$, which agree for all $X$ iff all products $\eta_{B,j}\eta_{A,j}$ are equal. With block-scalar
rates ($D_B=\eta_BI$, $D_A=\eta_AI$) this holds, the type (i) generators are the symmetric $X$ and the type (ii) ones the
antisymmetric $X$, and $\gl_r=\so_r\oplus\Sym_r$.

When all products equal $p$, $D_B=pD_A^{-1}$, and the type (i) generators are $X=SD_B$ with $S$ symmetric (then
$D_A^{-1}X=pD_A^{-1}SD_A^{-1}$ is symmetric too); the type (ii) generators are $X=KD_B$ with $K$ antisymmetric. With
$D=D_B^{1/2}$, conjugation $X\mapsto DXD^{-1}$ maps $SD_B$ to $DSD$ and $KD_B$ to $DKD$, so
$\mathfrak k\oplus\mathfrak p=D^{-1}(\so_r\oplus\Sym_r)D$. The isometries are the group elements $D^{-1}OD$ with
$O\in\SO(r)$, and $D^{-1}OD$ is orthogonal iff $O$ commutes with $D^2=D_B$. Unless $D_B$ is a multiple of the identity
(the block-scalar case), some of these isometries are not orthogonal, so they act on $Q$ by non-Euclidean maps and the
clipping caveat of (ii) applies.
\end{proof}

\begin{remark*}{instances and the discrete step}
Each instance listed after \thmref{Theorem}{III.5} applies (i) to the generator named there; we record the checks. All
assume gradient flow, or SGD without momentum, with per-block rates.
\begin{itemize}
\item \emph{LoRA.} For symmetric $X$, $J_{\xi_X}(B,A)=\tfrac12\big(\langle B,BX\rangle/\eta_B-\langle A,XA\rangle/\eta_A\big)=\tfrac12\tr(X\Phi)$
with $\Phi=B^\top B/\eta_B-AA^\top/\eta_A$, so the $r(r+1)/2$ charges are the entries of $\Phi$. The $\so_r$ generators
act by $(BO^\top,OA)$ with $O\in\SO(r)$, which are Euclidean isometries commuting with $\Lambda$, so (ii) holds under
global-norm clipping.
\item \emph{DyLoRA.} At each step the loss uses $\rho_b(B,A)=W_0+s_bBE_bA$ for a sampled $b$, with
$E_b=\diag(1,\dots,1,0,\dots,0)$ of rank $b$. The update is $-\Lambda D\rho_{b(t)}^\top\gamma$, and the proof of (i) needs
only $D\rho_{b(t)}(\xi q)=0$. Now $\rho_b(g\cdot q)=\rho_b(q)$ for all $(B,A)$ iff $g^{-1}E_bg=E_b$, that is, iff
$g$ preserves $\spanop(e_1,\dots,e_b)$ and $\spanop(e_{b+1},\dots,e_r)$. When every $b\in\{1,\dots,r\}$ can be sampled,
$g$ then preserves each line $\spanop(e_b)=\spanop(e_1,\dots,e_b)\cap\spanop(e_b,\dots,e_r)$, so it is diagonal. Diagonal $X$ are symmetric, which gives the charges $\diag\Phi$.
\item \emph{VeRA.} $\rho(b,d)=W_0+\diag(b)B_{\mathrm f}\diag(d)A_{\mathrm f}$ with frozen $B_{\mathrm f},A_{\mathrm f}$ is
invariant under $(b,d)\mapsto(cb,d/c)$, with generator $(b,-d)$ and charge $\tfrac12(\lVert b\rVert^2/\eta_b-\lVert d\rVert^2/\eta_d)$.
NOLA's two coefficient vectors behave in the same way.
\item \emph{AdaLoRA.} $\rho(P,\lambda,Q)=W_0+P\diag(\lambda)Q$ is invariant under $(p_k,\lambda_k)\mapsto(cp_k,\lambda_k/c)$ and
under $(\lambda_k,q_k)\mapsto(\lambda_k/c,cq_k)$; at equal rates the charges are proportional to
$\lVert p_k\rVert^2-\lambda_k^2$ and $\lVert q_k\rVert^2-\lambda_k^2$. The orthogonality regulariser
$\lVert P^\top P-I\rVert^2+\lVert QQ^\top-I\rVert^2$ is not invariant under these scalings, so the regularised loss does not
factor through $\rho$, and rank-budget masking changes $\rho$.
\item \emph{HRA.} $\rho(u_1,\dots,u_r)=W_0\prod_i(I-2u_iu_i^\top/\lVert u_i\rVert^2)$ is invariant under $u_i\mapsto cu_i$,
with charge $\tfrac12\lVert u_i\rVert^2/\eta$.
\item \emph{ReLU bottleneck adapters} $h\mapsto h+U\operatorname{relu}(Dh+\beta)$. Multiplying row $j$ of $D$ and $\beta_j$ by
$c>0$ and column $j$ of $U$ by $1/c$ preserves the function, since $\operatorname{relu}(cz)=c\operatorname{relu}(z)$. The
target is a function space, which \thmref{Observation}{III.2} allows, and the derivative of $\rho$ along the scaling
direction vanishes by homogeneity even at the kinks. Each hidden unit gives the charge
$\tfrac12\big((\lVert D_{j,:}\rVert^2+\beta_j^2)/\eta_D-\lVert U_{:,j}\rVert^2/\eta_U\big)$, with a common rate for $D$
and $\beta$.
\item \emph{TT adapters.} Neighbouring cores share a bond of dimension $r_k$, on which $\GL(r_k)$ acts by $h^{-1}$ on one
side and $h$ on the other. The symmetric generators give one conserved symmetric $r_k\times r_k$ matrix per bond, the
Gram matrix of the left unfolding of one core divided by its rate minus that of the right unfolding of the next core
divided by its rate.
\item \emph{DoRA's detached-norm backward, global-norm clipping, LoRA dropout.} DoRA writes $W=V\diag(\mathsf g/\lVert V\rVert_c)$
with $V=W_0+sBA$, column norms $\lVert V\rVert_c$ and magnitudes $\mathsf g$. Treating $\lVert V\rVert_c$ as a constant in
the backward pass gives $\nabla_B=s\gamma A^\top$ and $\nabla_A=sB^\top\gamma$ with
$\gamma=\nabla_WL\,\diag(\mathsf g/\lVert V\rVert_c)$. Clipping multiplies the factor gradients by one scalar, and dropout
replaces $G$ by the masked signal of \S\ref{proof:III.2}. In each case the update of $(B,A)$ is
$-\Lambda D\rho_{\mathrm{LoRA}}^\top\gamma$ for some $\gamma$, so $\Phi$ is conserved.
\item \emph{Discrete GD.} With $s$ absorbed into $G$, one step $B^+=B-\eta_BGA^\top$, $A^+=A-\eta_AB^\top G$ gives
\begin{gather*}
\tfrac{1}{\eta_B}B^{+\top}B^+=\tfrac{1}{\eta_B}B^\top B-(AG^\top B+B^\top GA^\top)+\eta_BAG^\top GA^\top,\\
\tfrac{1}{\eta_A}A^+A^{+\top}=\tfrac{1}{\eta_A}AA^\top-(B^\top GA^\top+AG^\top B)+\eta_AB^\top GG^\top B,
\end{gather*}
so $\Phi^+-\Phi=\eta_BAG^\top GA^\top-\eta_AB^\top GG^\top B$ exactly. This is of order $\eta$ while $\Phi$ is of order
$1/\eta$, so $\eta\Phi$, or $\Phi$ relative to its size, drifts by $O(\eta^2)$ per step.
\item \emph{Momentum and Adam.} A heavy-ball step adds a buffer of gradients taken at past points, $D\rho_{q_s}^\top G_s$
with $s<t$, and Adam rescales gradients coordinatewise. Neither step has the form $-\Lambda D\rho_q^\top\gamma$ at the
current $q$ with a constant $\Lambda$, so (i) does not apply.
\end{itemize}
\end{remark*}

% ====================================================================================================
\subsection{Corollary III.6}\label{proof:III.6}

\begin{restate}{Corollary}{III.6}{charge at initialisation; weight decay}
(a) $\Phi_0=B_0^\top B_0/\eta_B-A_0A_0^\top/\eta_A$ is the value of the conserved charge, fixed by the initialisation and
the rates. Since $\Ort(r)$ acts on $\Phi$ by conjugation and preserves $\theta$-trajectories (\thmref{Theorem}{III.5}(ii)),
only the spectrum of $\Phi_0$ labels a $\theta$-trajectory.

\smallskip
{\upshape\small\renewcommand{\arraystretch}{1.12}\noindent
\begin{tabularx}{\linewidth}{@{}>{\raggedright\arraybackslash}p{0.5\linewidth}>{\raggedright\arraybackslash}X@{}}
\toprule\headrow
\textbf{Initialisation} & \textbf{Charge $\Phi_0$}\\
\midrule
HF default LoRA on Linear and Conv layers ($B_0=0$; Kaiming-uniform $A_0$ with entries $U(-n^{-1/2},n^{-1/2})$, $n$ the fan-in, which is \texttt{in\_channels} times the kernel size for Conv) & $-A_0A_0^\top/\eta_A$, with $\mathbb{E}[A_0A_0^\top]=\tfrac13I_r$ and fluctuations $O(n^{-1/2})$ on and off the diagonal\\
HF default LoRA on embedding layers ($A_0=0$, normal $B_0$) & $+B_0^\top B_0/\eta_B$ (Init[B]-type)\\
HF \texttt{mica} \citep{rudiger2026mica} ($A_0=0$, frozen minor-singular $B_0$) & $+B_0^\top B_0/\eta_B$ (Init[B]-type, with $B$ frozen)\\
EVA \citep{paischer2024parameter} (orthonormal rows; default \texttt{whiten=False}) & $-I/\eta_A$ exactly\\
PiSSA in HF PEFT \citep{meng2024pissa}, and MiLoRA \citep{wang2024milora} if implemented like \texttt{pissa\_init} with the bottom-$r$ singular values (MiLoRA is not in HF PEFT); $B_0^\top B_0=A_0A_0^\top=S_r/s$, with $S_r$ the top-$r$, resp.\ bottom-$r$, singular values and $s$ the LoRA scale & $(S_r/s)(1/\eta_B-1/\eta_A)$, zero \textbf{iff} $\eta_A=\eta_B$ (when $S_r\ne0$)\\
OLoRA \citep{buyukakyuz2024olora} ($B_0=Q_r$, $A_0=R_r$) & $I/\eta_B-R_rR_r^\top/\eta_A$\\
``Init[B]'' \citep{hayou2024impact} ($A_0=0$) & $+B_0^\top B_0/\eta_B$\\
\bottomrule
\end{tabularx}}
\smallskip

(b) Let $G$ be any cotangent signal, with the scale $s$ absorbed into it. Under gradient flow, if decay acts per unit time
and identically on both factors,
\[\dot B=-\eta_BGA^\top-\lambda B,\qquad\dot A=-\eta_AB^\top G-\lambda A,\]
then $\dot\Phi=-2\lambda\Phi$. With learning-rate-coupled decay (PyTorch SGD \texttt{weight\_decay}; the decay term of
PyTorch AdamW), $\dot B=-\eta_B(GA^\top+\lambda B)$ and likewise for $A$, so that $\dot\Phi=-2\lambda(B^\top B-AA^\top)$.
Then $\Phi(t)=e^{-2\eta\lambda t}\Phi(0)$ holds when $\eta_A=\eta_B=\eta$. For $\eta_A\ne\eta_B$ it fails in general,
since $-2\lambda(B^\top B-AA^\top)$ is not a function of $\Phi$. With an actual Adam(W) gradient term, $\Phi$ drifts even
without decay, at first order in the step size; Adam lies outside \thmref{Theorem}{III.5}. So neither law describes AdamW
training.
\end{restate}

\begin{proof}
(a) \emph{Conservation and conjugation.} With block-scalar rates the entries of $\Phi$ are the charges of
\thmref{Theorem}{III.5}(i) (\S\ref{proof:III.5}, remark on instances), so under gradient flow $\Phi(t)=\Phi_0$, the value
fixed by $(B_0,A_0)$ and $(\eta_A,\eta_B)$. For $O\in\Ort(r)$, the map $k:(B,A)\mapsto(BO,O^\top A)$ is the gauge
element $g=O^\top$, and
\[\Phi(BO,O^\top A)=O^\top\tfrac{B^\top B}{\eta_B}O-O^\top\tfrac{AA^\top}{\eta_A}O=O^\top\Phi(B,A)\,O.\]
The linear map $k$ is orthogonal on $Q$ and commutes with block-scalar $\Lambda$, so $k\Lambda k^\top=\Lambda$, and
$\rho\circ k=\rho$. The trajectory argument of \thmref{Theorem}{III.5}(ii) uses only these two properties
(\S\ref{proof:III.5}), so the $\theta$-trajectories from $q_0$ and from $kq_0$ coincide for every signal that depends on
$q$ only through $(t,\rho(q))$, and also under global-norm clipping, since $k$ is a Euclidean isometry. This covers all of
$\Ort(r)$, reflections included. So initialisations whose charges are conjugate under $\Ort(r)$ produce the same
$\theta$-trajectory, and a function of $\Phi_0$ that labels $\theta$-trajectories must be invariant under
$\Phi_0\mapsto O^\top\Phi_0O$. By the spectral theorem the $\Ort(r)$-conjugacy class of the symmetric matrix $\Phi_0$ is
determined by its eigenvalues, so such a function is a function of the spectrum of $\Phi_0$.

\emph{The table.} Each entry follows by substituting the initialisation into $\Phi_0$.
\begin{itemize}
\item \emph{HF default, Linear and Conv.} $B_0=0$ gives $\Phi_0=-A_0A_0^\top/\eta_A$. HF initialises $A_0$ with
\texttt{kaiming\_uniform\_} at $a=\sqrt5$, whose bound is $\sqrt{2/(1+a^2)}\sqrt{3/n}=n^{-1/2}$, with $n$ the fan-in
(\texttt{in\_channels} times the product of the kernel sizes for Conv). An entry $U(-n^{-1/2},n^{-1/2})$ has variance
$1/(3n)$ and fourth moment $1/(5n^2)$. A diagonal entry of $A_0A_0^\top$ is a sum of $n$ independent squares, with mean
$1/3$ and variance $n\big(1/(5n^2)-1/(9n^2)\big)=4/(45n)$; an off-diagonal entry is a sum of $n$ independent products,
with mean $0$ and variance $n\cdot1/(9n^2)=1/(9n)$. So $\mathbb E[A_0A_0^\top]=\tfrac13I_r$, with fluctuations of standard
deviation $2/(3\sqrt{5n})$ on the diagonal and $1/(3\sqrt n)$ off it.
\item \emph{Embeddings, Init[B] and \texttt{mica}.} $A_0=0$ gives $\Phi_0=B_0^\top B_0/\eta_B$. For \texttt{mica}, $B$ is
frozen, so its rate is $0$ and \thmref{Theorem}{III.5} does not apply to it; the entry records the formula and its sign
type, and the training dynamics are those of the one-sided method $A\mapsto W_0+sB_0A$, along which $B_0^\top B_0$ is
constant and $AA^\top$ moves.
\item \emph{EVA.} $B_0=0$ and orthonormal rows, $A_0A_0^\top=I$, give $\Phi_0=-I/\eta_A$.
\item \emph{PiSSA in HF PEFT.} \texttt{pissa\_init} takes the top-$r$ singular triples $(U_r,S_r,V_r)$ of $W_0$, divides
$S_r$ by the scale $s$, and sets $A_0=(S_r/s)^{1/2}V_r^\top$ and $B_0=U_r(S_r/s)^{1/2}$. With $U_r^\top U_r=V_r^\top V_r=I$,
$B_0^\top B_0=A_0A_0^\top=S_r/s$, so $\Phi_0=(S_r/s)(1/\eta_B-1/\eta_A)$, which vanishes iff $\eta_A=\eta_B$ when
$S_r\ne0$. The same computation with the bottom-$r$ triples gives the MiLoRA entry.
\item \emph{OLoRA.} $B_0=Q_r$ has orthonormal columns, so $\Phi_0=I/\eta_B-R_rR_r^\top/\eta_A$.
\end{itemize}

(b) Write both decay conventions as $\dot B=-\eta_BGA^\top-c_BB$ and $\dot A=-\eta_AB^\top G-c_AA$, with $c_B=c_A=\lambda$
per unit time and $c_B=\eta_B\lambda$, $c_A=\eta_A\lambda$ when coupled. Then
\[\tfrac{d}{dt}\tfrac{B^\top B}{\eta_B}=-(AG^\top B+B^\top GA^\top)-2c_B\tfrac{B^\top B}{\eta_B},\qquad
\tfrac{d}{dt}\tfrac{AA^\top}{\eta_A}=-(B^\top GA^\top+AG^\top B)-2c_A\tfrac{AA^\top}{\eta_A},\]
and subtracting, the $G$-terms cancel:
\[\dot\Phi=-2c_B\tfrac{B^\top B}{\eta_B}+2c_A\tfrac{AA^\top}{\eta_A}.\]
Per unit time this is $\dot\Phi=-2\lambda\Phi$, so $\Phi(t)=e^{-2\lambda t}\Phi(0)$. Coupled, it is
$\dot\Phi=-2\lambda(B^\top B-AA^\top)$; when $\eta_A=\eta_B=\eta$ the right-hand side is $-2\lambda\eta\Phi$, so
$\Phi(t)=e^{-2\eta\lambda t}\Phi(0)$. For $\eta_A\ne\eta_B$ there are states with the same $\Phi$ and different
$B^\top B-AA^\top$: replacing $(B^\top B,AA^\top)$ by $(B^\top B+\eta_BC,AA^\top+\eta_AC)$ with $C\succeq0$ (realised by
suitable factors, since $m,n\ge r$) leaves $\Phi$ unchanged and changes $B^\top B-AA^\top$ by $(\eta_B-\eta_A)C$. Since
$\dot\Phi$ does not involve $G$, two such initial states have equal $\Phi(0)$ and different $\dot\Phi(0)$, so no law
$\dot\Phi=F(\Phi)$, and in particular no exponential law, holds for all trajectories. Some trajectories do satisfy one:
with $B=0$ and $G=0$, $\Phi(t)=e^{-2\eta_A\lambda t}\Phi(0)$; so the law fails in general and not always.

\emph{Adam.} One Adam step can move $\eta\Phi$ at first order. Take $m=n=r=1$, $s$ absorbed into $G$, rates
$\eta_B=\eta_A=\eta$, $(b,a)=(2,1)$, $G=1$ and zero optimizer state. The bias-corrected first step at $\varepsilon=0$ is
$\delta b=-\eta\operatorname{sign}(Ga)=-\eta$ and $\delta a=-\eta\operatorname{sign}(bG)=-\eta$, so
\[\eta\Phi^+-\eta\Phi=(2-\eta)^2-(1-\eta)^2-(4-1)=-2\eta,\]
of first order in $\eta$; the GD step of the previous remark moves $\eta\Phi$ by $O(\eta^2)$. Adam's step is not of the
form $-\Lambda D\rho^\top\gamma$, so it lies outside \thmref{Theorem}{III.5}, and neither decay law describes AdamW
training.
\end{proof}

\begin{remark*}{the two decay conventions as gauge fixings}
Per-time decay is the $\Lambda$-gradient flow of $L+\tfrac\lambda2\big(\lVert B\rVert^2/\eta_B+\lVert A\rVert^2/\eta_A\big)$,
since $-\eta_B\nabla_B\big(\tfrac\lambda2\lVert B\rVert^2/\eta_B\big)=-\lambda B$. Coupled decay is the $\Lambda$-gradient
flow of $L+\tfrac\lambda2(\lVert B\rVert^2+\lVert A\rVert^2)$. At an equilibrium of the coupled flow,
$\nabla_B=-\lambda B$ and $\nabla_A=-\lambda A$; since $B^\top\nabla_B=sB^\top GA^\top=\nabla_AA^\top$, this gives
$\lambda(B^\top B-AA^\top)=0$, so equilibria lie on the unweighted balanced slice $\mu=B^\top B-AA^\top=0$. At an
equilibrium of the per-time flow, $\eta_B\nabla_B=-\lambda B$ and $\eta_A\nabla_A=-\lambda A$, and the same identity gives
$\Phi=0$; per-time decay also drives $\Phi\to0$ exponentially by (b). The two slices
$\{B^\top B=AA^\top\}$ and $\{B^\top B=(\eta_B/\eta_A)AA^\top\}$ coincide iff $\eta_A=\eta_B$. HF's PiSSA has
$B_0^\top B_0=A_0A_0^\top$ at every rate, so it starts on the coupled-decay slice $\mu=0$, and off the per-time slice
$\Phi=0$ when $\eta_A\ne\eta_B$.
\end{remark*}

% ====================================================================================================
\subsection{Proposition III.7}\label{proof:III.7}

\begin{restate}{Proposition}{III.7}{Kempf--Ness slice; Loewner floor; known, \citet{xkempf1979length,xsrebro2004maximum}}
(a) On a full-rank $\GL_r$-orbit, the critical points of $\lVert B\rVert^2+\lVert A\rVert^2$ are exactly the balanced
representatives, those with $\mu=B^\top B-AA^\top=0$. They form one $\Ort(r)$-orbit, they are the minima, and the minimum
value is $2\lVert BA\rVert_*$.

(b) Assume gradient flow with block-scalar rates $\eta_A,\eta_B$, no weight decay, a loss that depends on $(B,A)$ only
through $BA$, and any cotangent signal of the form $\Lambda D\rho^\top\gamma$ (dropout, minibatches and clipping
included). From $B_0=0$,
\[B^\top B=\tfrac{\eta_B}{\eta_A}\big(AA^\top-A_0A_0^\top\big)\succeq0,\]
so $AA^\top\succeq A_0A_0^\top$ for all time; the $A$-factor never shrinks below its initialisation in the Loewner order.
$AA^\top(t)$ need not be monotone in $t$; only this floor is guaranteed. Discrete steps break the bound by $O(\eta^2)$ per
step. Under Adam it is not guaranteed and fails empirically. With per-time decay $\lambda$, only
$AA^\top\succeq e^{-2\lambda t}A_0A_0^\top$ holds.
\end{restate}

\begin{proof}
(a) Let $x=(B,A)$ have $\rk B=\rk A=r$, put $f(B,A)=\lVert B\rVert^2+\lVert A\rVert^2$, and let
$\mathcal O=\GL_r\cdot x$. Every factorisation $B'A'$ of $BA$ with inner dimension $r$ lies in $\mathcal O$: from
$\operatorname{col}B'=\operatorname{col}(BA)=\operatorname{col}B$ we get $B'=Bg^{-1}$ for a unique $g\in\GL_r$, and then
$Bg^{-1}A'=BA$ gives $A'=gA$ because $B$ is injective. As in Kempf--Ness theory (\bgref{D.30}, \auxref{B.26}), $x$ is a critical point
of $f$ on $\mathcal O$ when $e\in\GL_r$ is a critical point of $g\mapsto f(g\cdot x)$. With $e^{t\xi}\cdot x=(Be^{-t\xi},e^{t\xi}A)$,
\[\tfrac{d}{dt}\Big|_{t=0}f(e^{t\xi}\cdot x)=-2\langle B,B\xi\rangle+2\langle A,\xi A\rangle=-2\tr\big(\xi(B^\top B-AA^\top)\big)=-2\tr(\xi_s\mu),\]
where $\xi_s$ is the symmetric part of $\xi$ and the antisymmetric part drops out because $\mu$ is symmetric. This vanishes
for all $\xi$ iff $\mu=0$ (take $\xi=\mu$). So the critical points are exactly the balanced points of $\mathcal O$.

\emph{Convexity along symmetric directions.} The function $f$ is $\Ort(r)$-invariant, $f(BO^\top,OA)=f(B,A)$, and $\mu$
is $\Ort(r)$-equivariant, $\mu(O\cdot x)=O\mu(x)O^\top$. By the polar decomposition every $g\in\GL_r$ is $g=Oe^\xi$ with
$O\in\Ort(r)$ and $\xi$ symmetric (\bgref{D.27}, \auxref{B.19}(a); \citealp[\S7.3]{xhorn2013matrix}), so every point of $\mathcal O$ is
$O\cdot(e^\xi\cdot x)$ and $f(g\cdot x)=f(e^\xi\cdot x)$. For symmetric $\xi=R\diag(d)R^\top$ with $R$ orthogonal, put
$h(t)=f(e^{t\xi}\cdot x)$. Orthogonal invariance of the Frobenius norm gives
\[h(t)=\lVert BRe^{-t\diag(d)}\rVert^2+\lVert e^{t\diag(d)}R^\top A\rVert^2=\sum_i\Big(\lVert(BR)_{:,i}\rVert^2e^{-2td_i}+\lVert(R^\top A)_{i,:}\rVert^2e^{2td_i}\Big),\]
\[h''(t)=4\sum_id_i^2\Big(\lVert(BR)_{:,i}\rVert^2e^{-2td_i}+\lVert(R^\top A)_{i,:}\rVert^2e^{2td_i}\Big)\ge0 .\]
So $h$ is convex.

\emph{Balanced points are the minima.} Let $x$ be balanced and $y=O\cdot(e^\xi\cdot x)\in\mathcal O$. Then
$h'(0)=-2\tr(\xi\mu(x))=0$, and convexity gives $f(y)=h(1)\ge h(0)=f(x)$. So every balanced point is a global minimum on
$\mathcal O$; conversely every minimum is a critical point, hence balanced.

\emph{One $\Ort(r)$-orbit.} Let $x$ and $y=O\cdot(e^\xi\cdot x)$ both be balanced. Then $e^\xi\cdot x=O^{-1}\cdot y$ is
balanced by equivariance of $\mu$, so $h'(0)=0$ and $h'(1)=-2\tr\big(\xi\mu(e^\xi\cdot x)\big)=0$. Since $h'$ is
non-decreasing, $h'\equiv0$ on $[0,1]$, so $h''(0)=0$. Every column $(BR)_{:,i}$ is nonzero because $BR$ has full column
rank, so all $d_i=0$, $\xi=0$ and $y=O\cdot x$. Conversely the $\Ort(r)$-orbit of a balanced point consists of balanced
points.

\emph{Existence and value.} Let $BA=U\Sigma V^\top$ be a compact SVD (\auxref{B.11}(a)) with $\Sigma\in\R^{r\times r}$ positive diagonal. The
factorisation $(U\Sigma^{1/2},\Sigma^{1/2}V^\top)$ lies in $\mathcal O$ by the first paragraph, is balanced since both
Gram matrices equal $\Sigma$, and has $f=2\tr\Sigma=2\lVert BA\rVert_*$. This is the minimum value.

(b) Under these hypotheses the update has the form of \thmref{Theorem}{III.5}, so $\Phi$ is conserved. With $B_0=0$,
$\Phi_0=-A_0A_0^\top/\eta_A$, and $B^\top B/\eta_B-AA^\top/\eta_A=-A_0A_0^\top/\eta_A$ rearranges to
$B^\top B=\tfrac{\eta_B}{\eta_A}(AA^\top-A_0A_0^\top)$. Since $B^\top B\succeq0$, $AA^\top\succeq A_0A_0^\top$.

\emph{No monotonicity.} Take $m=n=r=1$, $a(0)=a_0>0$, $b(0)=0$ and the signal $\gamma=-1$ on $[0,1)$ and $\gamma=+1$ on
$[1,2]$, the gradients of the minibatch losses $-w$ and $w$. The flow is $\dot b=-\eta_Bs\gamma a$, $\dot a=-\eta_As\gamma b$.
On $[0,1]$, $a(t)=a_0\cosh(s\sqrt{\eta_A\eta_B}\,t)$ increases. On $[1,2]$ the system is the time reversal of the first one,
so it retraces its path and returns to $(a_0,0)$ at $t=2$, and $a$ decreases. The floor $a^2\ge a_0^2$ holds throughout,
but $a^2$ is not monotone.

\emph{Discrete steps.} By the discrete formula of \S\ref{proof:III.5}, after $k$ GD steps (with $s$ absorbed into $G_j$)
\begin{align*}
A_kA_k^\top&=\eta_A\Big(\tfrac{1}{\eta_B}B_k^\top B_k-\Phi_k\Big)\succeq-\eta_A\Phi_k
=A_0A_0^\top-\sum_{j<k}\big(\eta_A\eta_BA_jG_j^\top G_jA_j^\top-\eta_A^2B_j^\top G_jG_j^\top B_j\big)\\
&\succeq A_0A_0^\top-\eta_A\eta_B\sum_{j<k}A_jG_j^\top G_jA_j^\top,
\end{align*}
so the floor holds up to an error of order $\eta^2$ per step. The identity itself breaks at that order already at the first
step, where $B_1=-\eta_BG_0A_0^\top$ and $A_1=A_0$ give $B_1^\top B_1=\eta_B^2A_0G_0^\top G_0A_0^\top$ against
$\tfrac{\eta_B}{\eta_A}(A_1A_1^\top-A_0A_0^\top)=0$.

\emph{Adam.} Take $m=n=r=1$, $s=1$, $(b_0,a_0)=(0,1)$, minibatch gradients $G_1=1$ and $G_2=-1$, any $\varepsilon\ge0$
and bias correction. At step 1, $\nabla_a=b_0G_1=0$, so $a_1=1$, and $b_1=-\eta_B\phi(G_1a_0)$ has sign $-1$. At step 2,
$\nabla_a=b_1G_2>0$ is the first nonzero gradient of $a$, so Adam's step on $a$ is
$-\eta_A\hat m/(\sqrt{\hat v}+\varepsilon)$ with $\hat m=\nabla_a/(1+\beta_1)>0$; at $\varepsilon=0$ it equals
$-\eta_A\sqrt{1+\beta_2}/(1+\beta_1)$. So $a_2<1=a_0$ for every $\varepsilon\ge0$, and at $\varepsilon=0$ the violation is of first order in $\eta_A$.

\emph{Per-time decay.} By \thmref{Corollary}{III.6}(b), $\Phi(t)=e^{-2\lambda t}\Phi(0)=-e^{-2\lambda t}A_0A_0^\top/\eta_A$,
so $AA^\top=e^{-2\lambda t}A_0A_0^\top+\tfrac{\eta_A}{\eta_B}B^\top B\succeq e^{-2\lambda t}A_0A_0^\top$.
\end{proof}

\begin{remark*}{nuclear-norm penalties}
By (a), $\min\{\lambda(\lVert B\rVert^2+\lVert A\rVert^2):sBA=\Delta W\}=2\lambda\lVert\Delta W\rVert_*/s$ for $\Delta W$ of
rank $r$; \citet{xsrebro2004maximum} show (\auxref{B.21}) that $\min\{\tfrac12(\lVert B\rVert^2+\lVert A\rVert^2):BA=Z\}=\lVert Z\rVert_*$
for every inner dimension at least $\rk Z$, which covers rank-deficient $\Delta W$. So the explicit penalty
$\lambda(\lVert B\rVert^2+\lVert A\rVert^2)$ matches a $(2\lambda/s)$-nuclear-norm penalty at its minimisers, and coupled
decay with coefficient $\lambda$, the flow of $L+\tfrac\lambda2(\lVert B\rVert^2+\lVert A\rVert^2)$, matches a
$(\lambda/s)$-nuclear-norm penalty.
\end{remark*}

% ====================================================================================================
\subsection{Proposition III.8}\label{proof:III.8}

\begin{restate}{Proposition}{III.8}{the balance--escape trilemma}
No LoRA initialisation is simultaneously
\begin{enumerate}[label=(\alph*)]
\item identity-at-init without a residual split ($B_0A_0=0$),
\item balanced ($B_0^\top B_0=cA_0A_0^\top$ for some $c>0$: unweighted, $c=1$, or rate-weighted, $\Phi_0=0$ with
$c=\eta_B/\eta_A$), and
\item non-stationary, meaning that $D\rho$ at $(B_0,A_0)$ is not the zero map; equivalently, $\dot W(0)\ne0$ for some
parameter velocity, under any first-order optimizer.
\end{enumerate}
Moreover, if $B_0A_0=0$ then $\langle B_0^\top B_0,A_0A_0^\top\rangle_F=0$, so
\[\lVert B_0^\top B_0-cA_0A_0^\top\rVert_F^2=\lVert B_0^\top B_0\rVert_F^2+c^2\lVert A_0A_0^\top\rVert_F^2,\]
and zero-product initialisations are maximally unbalanced.
\end{restate}

\begin{proof}
Assume (a) and (b). Then $0=B_0^\top(B_0A_0)=(B_0^\top B_0)A_0=cA_0A_0^\top A_0$, and multiplying by $A_0^\top$ on the left
gives $c(A_0^\top A_0)^2=0$. The matrix $A_0^\top A_0$ is symmetric positive semidefinite with vanishing square, so its
eigenvalues vanish and $A_0^\top A_0=0$, that is, $A_0=0$. Then $B_0^\top B_0=cA_0A_0^\top=0$, so $B_0=0$. At $(0,0)$,
$D\rho(\dot B,\dot A)=s(\dot B\cdot0+0\cdot\dot A)=0$, which contradicts (c). The equivalence in (c) is the definition of
the zero map: some parameter velocity has $\dot W=D\rho(\dot B,\dot A)\ne0$ iff $D\rho\ne0$. At $(0,0)$ also
$\nabla_B=sGA^\top=0$ and $\nabla_A=sB^\top G=0$, so every update rule built from these gradients and the current
parameters, including Adam, momentum and weight decay started from zero state, stays at $(0,0)$.

For the identity, $\langle B_0^\top B_0,A_0A_0^\top\rangle_F=\tr(B_0^\top B_0A_0A_0^\top)=\tr\big(B_0^\top(B_0A_0)A_0^\top\big)=0$,
and expanding the square gives the displayed formula. For positive semidefinite $P$ and $N$,
$\langle P,N\rangle_F=\tr(P^{1/2}NP^{1/2})\ge0$, so $\lVert P-cN\rVert_F^2\le\lVert P\rVert_F^2+c^2\lVert N\rVert_F^2$, with
equality iff $\langle P,N\rangle_F=0$. Zero-product initialisations attain this maximum for the given norms of the two
Gram matrices.
\end{proof}

% ====================================================================================================
\subsection{Theorem III.9}\label{proof:III.9}

\begin{restate}{Theorem}{III.9}{closed-form isotropic-charge LoRA dynamics; known, \citet{xtarmoun2021understanding} and \citet{xmin2021explicit}}
Let $\rho(B,A)=W_{\mathrm{res}}+sBA$ be trained by gradient flow with rates $\eta_A,\eta_B$, and suppose the charge is
isotropic, $\Phi=-\kappa I_r$. Put $\mu=s\sqrt{\eta_A\eta_B}$ (a rate, unrelated to the moment map $\mu$ of
\thmref{Proposition}{III.7}), and on the locus where $Z=W-W_{\mathrm{res}}$ has rank $r$ write its compact SVD
$Z=U\Sigma V^\top$. Then
\[\dot W=-\mu^2\big[G\,V g_\kappa(\Sigma^2/\mu^2)V^\top+U f_\kappa(\Sigma^2/\mu^2)U^\top G\big],\qquad
f_\kappa(x)=\sqrt{x+\kappa^2/4}-\kappa/2,\quad g_\kappa=f_\kappa+\kappa .\]
The identity holds pointwise on $\{\Phi=-\kappa I,\ \rk Z=r\}$; gradient flow is used only to keep $\Phi$ fixed. The two
terms under the root are equal at the \emph{crossover scale} $\sigma^*=\mu\lvert\kappa\rvert/2$.
\begin{itemize}
\item For $\sigma\ll\sigma^*$ with $\kappa>0$, $f\to0$ and $g\to\kappa$, so $\dot W\approx-\mu^2\kappa\,GP_V$. These are
LoRA-FA dynamics. The contribution of $A$ to $\dot W$ is $O(\sigma^2/(\mu^2\kappa))$, although $A$ itself still drifts
slowly. For $\kappa<0$ the roles swap, $\dot W\approx-\mu^2\lvert\kappa\rvert P_UG$, and only $A$ effectively moves
(Init[B]).
\item For $\sigma\gg\sigma^*$, $f\approx g\approx\sigma/\mu$, so $\dot W\approx-\mu\big[G(Z^\top Z)^{1/2}+(ZZ^\top)^{1/2}G\big]$.
These are the balanced dynamics of Arora, Cohen and Hazan.
\item \emph{Off the rank-$r$ locus.} Write $\dot W=-\mu^2(GS+TG)$ with $S=A^\top A/\eta_A$ and $T=BB^\top/\eta_B$, so that
on the rank-$r$ locus $S=Vg_\kappa(\Sigma^2/\mu^2)V^\top$ and $T=Uf_\kappa(\Sigma^2/\mu^2)U^\top$. For $\kappa\ge0$,
$T=f_\kappa(ZZ^\top/\mu^2)$ still holds as a full matrix function, but $S$ gains the term
$\kappa(P_{\operatorname{row}A}-P_V)$; symmetrically, for $\kappa<0$, $T$ gains $\lvert\kappa\rvert(P_{\operatorname{col}B}-P_U)$.
At the default-LoRA start $Z=0$ ($B_0=0$) this leaves $S=\kappa P_{\operatorname{row}A_0}$ and $T=0$, so
$\dot W=-\mu^2\kappa\,GP_{\operatorname{row}A_0}$, and the $\sigma\ll\sigma^*$ regime extends continuously to $t=0$.
\end{itemize}
\end{restate}

\begin{proof}
Let $G$ be the signal driving the flow, $\dot B=-\eta_BsGA^\top$ and $\dot A=-\eta_AsB^\top G$; under gradient flow
$G=\nabla_\theta L$, and \thmref{Theorem}{III.5}(i) keeps $\Phi=-\kappa I$ for all time if it holds at $t=0$. Everything
below uses only the value of $\Phi$ at the point considered.

\emph{Normalisation.} Put $\beta=B/\sqrt{\eta_B}$ and $\alpha=A/\sqrt{\eta_A}$. Then
$\beta^\top\beta-\alpha\alpha^\top=\Phi=-\kappa I$ and $\beta\alpha=BA/\sqrt{\eta_A\eta_B}=Z/\mu$. Let
$T=\beta\beta^\top=BB^\top/\eta_B$ and $S=\alpha^\top\alpha=A^\top A/\eta_A$. Then
\[\dot W=s(\dot BA+B\dot A)=-s^2\big(\eta_BGA^\top A+\eta_ABB^\top G\big)=-s^2\eta_A\eta_B(GS+TG)=-\mu^2(GS+TG).\]

\emph{Two quadratic identities.} Using $\alpha\alpha^\top=\beta^\top\beta+\kappa I$ and $\beta^\top\beta=\alpha\alpha^\top-\kappa I$,
\[\tfrac{ZZ^\top}{\mu^2}=\beta\alpha\alpha^\top\beta^\top=\beta(\beta^\top\beta+\kappa I)\beta^\top=T^2+\kappa T,\qquad
\tfrac{Z^\top Z}{\mu^2}=\alpha^\top\beta^\top\beta\alpha=\alpha^\top(\alpha\alpha^\top-\kappa I)\alpha=S^2-\kappa S .\]
These hold at every point with $\Phi=-\kappa I$.

\emph{Spectra on the rank-$r$ locus.} If $\rk Z=r$, then $\beta$ and $\alpha^\top$ have full column rank, so
$\beta^\top\beta\succ0$ and $\alpha\alpha^\top\succ0$. The nonzero eigenvalues of $T=\beta\beta^\top$ are the eigenvalues of
$\beta^\top\beta=\alpha\alpha^\top-\kappa I$ (\auxref{B.36}(c)), which are positive and exceed $-\kappa$; so they exceed $\max(0,-\kappa)$.
Likewise the nonzero eigenvalues of $S$ are those of $\alpha\alpha^\top=\beta^\top\beta+\kappa I$, which exceed $\max(0,\kappa)$.
The range of $T$ is $\operatorname{col}\beta=\operatorname{col}Z=U$, and that of $S$ is $\operatorname{row}\alpha=\operatorname{row}Z=V$.

\emph{Inversion.} The polynomial $p(t)=t^2+\kappa t$ is strictly increasing on $[-\kappa/2,\infty)$, which contains the
nonzero eigenvalues of $T$ because $\max(0,-\kappa)\ge-\kappa/2$, and $p(t)>0$ there while $p(0)=0$. So $p$ is injective
on the spectrum of $T$, the eigenspaces of $T$ are those of $p(T)=ZZ^\top/\mu^2=U(\Sigma^2/\mu^2)U^\top$ (by the compact SVD of $Z$,
\auxref{B.11}(a)), and on the
eigenvector $u_i$ with $ZZ^\top u_i=\sigma_i^2u_i$ the eigenvalue of $T$ is the unique root $t>\max(0,-\kappa)$ of
$t^2+\kappa t=\sigma_i^2/\mu^2$, namely $f_\kappa(\sigma_i^2/\mu^2)$. On $U^\perp$ both $T$ and $ZZ^\top$ vanish. Hence
$T=Uf_\kappa(\Sigma^2/\mu^2)U^\top$. In the same way $q(s)=s^2-\kappa s$ is strictly increasing on $[\kappa/2,\infty)$,
the root above $\max(0,\kappa)$ of $s^2-\kappa s=x$ is $\kappa/2+\sqrt{x+\kappa^2/4}=g_\kappa(x)$, and
$S=Vg_\kappa(\Sigma^2/\mu^2)V^\top$. The compact SVD matters: $f_\kappa(0)=\lvert\kappa\rvert$ for $\kappa<0$ and
$g_\kappa(0)=\kappa$ for $\kappa>0$, so the full matrix functions $f_\kappa(ZZ^\top/\mu^2)$ and
$g_\kappa(Z^\top Z/\mu^2)$ would put nonzero values on $\ker Z^\top$ and $\ker Z$, where $T$ and $S$ vanish. Substituting
into $\dot W=-\mu^2(GS+TG)$ gives the formula. The crossover: $x=\sigma^2/\mu^2$ equals $\kappa^2/4$ iff
$\sigma=\mu\lvert\kappa\rvert/2$.

\emph{The regimes.} For $\kappa>0$, $f_\kappa(x)=x/\big(\sqrt{x+\kappa^2/4}+\kappa/2\big)\in[0,x/\kappa]$ and
$g_\kappa(x)-\kappa=f_\kappa(x)$. So $\dot W=-\mu^2\kappa\,GP_V+E$ with $E=-\mu^2\big(GVf_\kappa(\Sigma^2/\mu^2)V^\top+TG\big)$
and $\lVert E\rVert\le2(\sigma_1^2/\kappa)\lVert G\rVert$, where $\sigma_1$ is the largest singular value of $Z$. The
contribution $-\mu^2TG=s\,B\dot A$ of $A$ has gain at most $f_\kappa(\sigma_1^2/\mu^2)\le\sigma_1^2/(\mu^2\kappa)$, against
the gain $g_\kappa\ge\kappa$ of the $B$-term. The factor $A$ still moves, at speed
$\lVert\dot A\rVert\le\eta_As\lVert B\rVert_2\lVert G\rVert$ with $\lVert B\rVert_2^2=\eta_B\lVert T\rVert_2\le\eta_B\sigma_1^2/(\mu^2\kappa)$.
For $\kappa<0$ the same bounds hold with the roles exchanged: $g_\kappa(x)\in[0,x/\lvert\kappa\rvert]$ and
$f_\kappa=g_\kappa+\lvert\kappa\rvert$, so $\dot W=-\mu^2\lvert\kappa\rvert P_UG+O(\sigma_1^2/\lvert\kappa\rvert)\lVert G\rVert$.
For large $x$, $0\le\sqrt{x+\kappa^2/4}-\sqrt x\le\lvert\kappa\rvert/2$ gives $\lvert f_\kappa(x)-\sqrt x\rvert\le\lvert\kappa\rvert$
and $\lvert g_\kappa(x)-\sqrt x\rvert\le\lvert\kappa\rvert$. At $x=\sigma_i^2/\mu^2$ the relative error is at most
$\lvert\kappa\rvert\mu/\sigma_i=2\sigma^*/\sigma_i$, so for $\sigma_i\gg\sigma^*$, $f\approx g\approx\sigma_i/\mu$ and
\[\dot W\approx-\mu^2\big[GV(\Sigma/\mu)V^\top+U(\Sigma/\mu)U^\top G\big]=-\mu\big[G(Z^\top Z)^{1/2}+(ZZ^\top)^{1/2}G\big].\]
This is the end-to-end dynamics of a balanced depth-2 linear network with rate $\mu$ \citep[Thm.~1 with
$N=2$]{arora2018optimization}; at $\kappa=0$ it is exact, since $f_0(x)=g_0(x)=\sqrt x$.

\emph{Off the rank-$r$ locus.} Let $\kappa>0$ and $\rk Z<r$. The nonzero eigenvalues of $T$ are those of
$\beta^\top\beta\succeq0$, and on them $p(t)=t^2+\kappa t>0$; if $ZZ^\top v=0$ then $T(T+\kappa I)v=0$ with $T+\kappa I\succ0$,
so $Tv=0$. Hence $\ker T=\ker ZZ^\top$ and the inversion above gives $T=f_\kappa(ZZ^\top/\mu^2)$ as a full matrix function,
using $f_\kappa(0)=0$ for $\kappa\ge0$. For $S$, decompose $\operatorname{row}A=\operatorname{row}Z\oplus(\operatorname{row}A\cap\ker Z)$
orthogonally, which is possible since $\operatorname{row}Z\subseteq\operatorname{row}A$ and $(\operatorname{row}Z)^\perp=\ker Z$.
On $(\operatorname{row}A)^\perp=\ker\alpha$, $S=0$. If $v=\alpha^\top w\in\operatorname{row}A$ has $Zv=0$, then
$0=\beta\alpha\alpha^\top w=\beta(\beta^\top\beta+\kappa I)w$, so $\beta^\top\beta(\beta^\top\beta+\kappa I)w=0$; as
$\beta^\top\beta+\kappa I\succ0$ commutes with $\beta^\top\beta\succeq0$, this forces $\beta^\top\beta w=0$, hence $\beta w=0$,
and $Sv=\alpha^\top(\alpha\alpha^\top)w=\alpha^\top(\beta^\top\beta+\kappa I)w=\kappa v$. On $\operatorname{row}Z$ the
eigenvalues of $S$ exceed $\kappa$ and the inversion gives $Vg_\kappa(\Sigma^2/\mu^2)V^\top$. Together,
$S=Vg_\kappa(\Sigma^2/\mu^2)V^\top+\kappa(P_{\operatorname{row}A}-P_V)$. The case $\kappa<0$ is symmetric under
$(B,A,Z,\eta_B,\eta_A,\kappa)\mapsto(A^\top,B^\top,Z^\top,\eta_A,\eta_B,-\kappa)$, which exchanges $S$ and $T$. At $\kappa=0$ the correction terms vanish. At $Z=0$ with $B_0=0$
and $A_0A_0^\top=cI$, $\kappa=c/\eta_A>0$, $T=0$ and $S=\kappa P_{\operatorname{row}A_0}$, so
$\dot W=-\mu^2\kappa\,GP_{\operatorname{row}A_0}=-s^2\eta_BcGP_{\operatorname{row}A_0}$, which agrees with the direct value
$-s^2\eta_BGA_0^\top A_0$ because $A_0^\top A_0=cP_{\operatorname{row}A_0}$. As $Z\to0$ along the flow with $V\to\operatorname{row}A_0$,
the rank-$r$ formula tends to the same value, since $g_\kappa(0)=\kappa$ and $f_\kappa(0)=0$.
\end{proof}

% ====================================================================================================
\subsection{Theorem III.10}\label{proof:III.10}

\begin{restate}{Theorem}{III.10}{the hyperparameter torus; the LoRA/Adam case is known, \citet{schulman2025lora}}
Let $\rho(q)=\theta_b+s\,\delta(q_1,\dots,q_k)$ with $\delta$ multihomogeneous of multidegree $(d_1,\dots,d_k)$. Train block
$i$ from $q_i(0)=\sigma_iu_i$ with Adam ($\beta_1\in[0,1)$, $\beta_2\in(0,1)$, bias correction allowed) at learning rate
$\eta_i(t)$, on the loss $L\circ\rho$ with no $q$-dependent regulariser. Assume no gradient clipping (or per-block
clipping with thresholds rescaled as below), common random numbers (the same data order and dropout masks in both runs),
and exact arithmetic (the result is bit-exact in floating point when every $c_i$ is a power of $2$). Then the trajectory
$t\mapsto\rho(q(t))$ is invariant under the torus action of $c\in\R_{>0}^k$:
\[\big(s,\ \eta_i(t),\ \sigma_i,\ \varepsilon_i,\ \lambda_i,\ C_i\big)\mapsto\Big(s\prod_ic_i^{-d_i},\ c_i\eta_i(t),\ c_i\sigma_i,\ \varepsilon_i/c_i,\ \lambda_i/c_i,\ C_i/c_i\Big).\]
Here $\varepsilon_i$ is a per-block Adam $\varepsilon$ (with the convention $0/0=0$ at $\varepsilon=0$), $\lambda_i$ a
per-block decoupled (AdamW) weight decay, and $C_i$ a per-block clipping threshold. Under SGD, with or without heavy-ball
or Nesterov momentum, the same holds with $\eta_i(t)\mapsto c_i^2\eta_i(t)$ and coupled $L^2$ decay
$\lambda_i\mapsto\lambda_i/c_i^2$. In particular, any learning-rate schedule shared by all blocks, warm-up included, is
preserved.
\end{restate}

\begin{proof}
\emph{The rescaling is an isomorphism.} Write $c\odot q=(c_1q_1,\dots,c_kq_k)$ and $\rho'=\theta_b+s'\delta$ with
$s'=s\prod_ic_i^{-d_i}$. Multihomogeneity (\bgref{D.19}) gives
$\rho'(c\odot q)=\theta_b+s\prod_ic_i^{-d_i}\prod_ic_i^{d_i}\delta(q)=\rho(q)$ for all $q$, so $q\mapsto c\odot q$ is a
diffeomorphism over $\Theta$ that maps the base point $q(0)$ to $q'(0)=c\odot q(0)$; it is an isomorphism in $\Meth$. Let
$\ell_t=L_t\circ\rho$ and $\ell'_t=L_t\circ\rho'$ be the step-$t$ losses, with the same minibatch and dropout masks in
both runs and no $q$-dependent regulariser. Differentiating $\ell'_t(c\odot q)=\ell_t(q)$ gives
\begin{equation}\label{eq:C2-torus-grad}
\nabla_{q_i}\ell'_t(c\odot q)=c_i^{-1}\nabla_{q_i}\ell_t(q).
\end{equation}

\emph{The optimizer.} For block $i$ at step $t\ge1$, Adam(W) computes the gradient $g_t$ of $\ell_t$ at $q_{t-1}$, clips it
per block, $g_t\leftarrow g_t\min(1,C_i/\lVert g_t\rVert)$ (with factor $1$ when $g_t=0$, and no clipping step when there
is no clipping), and sets
\begin{gather*}
m_t=\beta_1m_{t-1}+(1-\beta_1)g_t,\qquad v_t=\beta_2v_{t-1}+(1-\beta_2)\,g_t\odot g_t,\\
q_t=q_{t-1}-\eta_i(t)\Big(\frac{m_t/b_{1,t}}{\sqrt{v_t/b_{2,t}}+\varepsilon_i}+\lambda_iq_{t-1}\Big),
\end{gather*}
with $m_0=v_0=0$, entrywise square root and division, $b_{j,t}=1-\beta_j^t$ with bias correction and $b_{j,t}=1$ without,
and $0/0=0$. PyTorch's AdamW first multiplies $q_{t-1}$ by $1-\eta_i(t)\lambda_i$ and then takes the Adam step, which is the
same update.

\emph{Induction.} Run the original hyperparameters from $q(0)$ and the transformed ones from $q'(0)=c\odot q(0)$, which is
what $\sigma_i\mapsto c_i\sigma_i$ means. Suppose that after step $t-1$, in every block,
\[q'_{t-1}=c_iq_{t-1},\qquad m'_{t-1}=c_i^{-1}m_{t-1},\qquad v'_{t-1}=c_i^{-2}v_{t-1},\]
which holds at $t=1$. By \eqref{eq:C2-torus-grad} the raw gradients satisfy $g'_t=c_i^{-1}g_t$. With $C'_i=C_i/c_i$ the
clipping factors agree, $\min(1,C'_i/\lVert g'_t\rVert)=\min(1,C_i/\lVert g_t\rVert)$, so the clipped gradients also satisfy
$g'_t=c_i^{-1}g_t$. The moment recursions are linear in $g_t$ and $g_t\odot g_t$, so $m'_t=c_i^{-1}m_t$ and
$v'_t=c_i^{-2}v_t$. With $\varepsilon'_i=\varepsilon_i/c_i$ the normalised step is unchanged,
\[\frac{m'_t/b_{1,t}}{\sqrt{v'_t/b_{2,t}}+\varepsilon'_i}=\frac{c_i^{-1}m_t/b_{1,t}}{c_i^{-1}\big(\sqrt{v_t/b_{2,t}}+\varepsilon_i\big)}=\frac{m_t/b_{1,t}}{\sqrt{v_t/b_{2,t}}+\varepsilon_i},\]
including the entries where $v_t=0$, at which $m_t=0$ as well and both sides are $0$. The decay term transforms as
$\lambda'_iq'_{t-1}=(\lambda_i/c_i)c_iq_{t-1}=\lambda_iq_{t-1}$. With $\eta'_i(t)=c_i\eta_i(t)$ the whole step is multiplied
by $c_i$, so $q'_t=c_iq_t$, which closes the induction. Then $\rho'(q'_t)=\rho(q_t)$ for all $t$. The value of $\eta_i(t)$
at each step was arbitrary, so any schedule works; a schedule shared by all blocks, $\eta_i(t)=\eta_i\varphi(t)$, becomes
$c_i\eta_i\varphi(t)$, the same schedule.

\emph{SGD.} With coupled $L^2$ decay the step-$t$ direction is $g_t+\lambda_iq_{t-1}$, the heavy-ball buffer is
$b_t=\mu_mb_{t-1}+(1-\tau)(g_t+\lambda_iq_{t-1})$ (with momentum $\mu_m$ and dampening $\tau$), the direction is $b_t$, or
$g_t+\lambda_iq_{t-1}+\mu_mb_t$ for Nesterov, and $q_t=q_{t-1}-\eta_i(t)\cdot(\text{direction})$. Suppose $q'_{t-1}=c_iq_{t-1}$
and $b'_{t-1}=c_i^{-1}b_{t-1}$. Then $g'_t=c_i^{-1}g_t$ and, with $\lambda'_i=\lambda_i/c_i^2$,
$\lambda'_iq'_{t-1}=c_i^{-1}\lambda_iq_{t-1}$; every quantity entering the direction is multiplied by $c_i^{-1}$, so
$b'_t=c_i^{-1}b_t$ and the direction is multiplied by $c_i^{-1}$. With $\eta'_i(t)=c_i^2\eta_i(t)$ the step is multiplied by
$c_i$, and $q'_t=c_iq_t$. Per-block clipping with $C_i/c_i$ is handled as for Adam.

\emph{Floating point.} If every $c_i$ is a power of $2$, multiplication by $c_i^{\pm1}$ and by $\prod_ic_i^{\mp d_i}$ is
exact, and it commutes with every correctly rounded addition, multiplication, division and square root, barring overflow
and underflow (\auxref{B.39}). Then the forward pass computes bit-identical weights and outputs in both
runs, the backward pass computes gradients that are exactly $c_i^{-1}$ times the original ones, and every quantity in the
induction above holds exactly in floating point, so the two trajectories agree bit for bit.
\end{proof}

% ====================================================================================================
\subsection{Corollary III.11}\label{proof:III.11}

\begin{restate}{Corollary}{III.11}{LoRA+ is $\alpha$ in disguise; rsLoRA is LoRA; (a) and the LoRA+ reading are known, \citet{schulman2025lora}}
For linear LoRA modules ($k=2$, $d=(1,1)$, $\sigma_B=0$), under the hypotheses of \thmref{Theorem}{III.10}:

(a) \emph{Effective hyperparameters.} Under Adam the torus invariants of $(s,\eta_A,\eta_B,\sigma_A)$ are
$(s\eta_A\eta_B,\ \sigma_A/\eta_A)$; the quantities $\alpha\cdot\mathrm{init}_A\cdot\mathrm{LR}_B$ and
$\mathrm{init}_A/\mathrm{LR}_A$ of Schulman et al.\ are equivalent generators. Under SGD the invariants are
$(s^2\eta_A\eta_B,\ \sigma_A^2/\eta_A)$.

(b) \emph{LoRA+.} Under Adam with $\varepsilon=0$, no weight decay, no gradient clipping and a learning-rate schedule
shared by both groups, LoRA+ with ratio $\lambda$ at scale $s$ produces \emph{exactly} the trajectory of plain LoRA at
scale $\lambda s$ (same $\eta_A$, same $A_0$). For $\varepsilon>0$ the equivalence needs a per-block
$\varepsilon_B'=\lambda\varepsilon$, and with AdamW a per-group decay $\lambda\cdot\mathrm{wd}$ on $B$. Under SGD it holds at
scale $\sqrt\lambda\,s$, or equivalently with plain LoRA at the single rate $\sqrt{\eta_A\eta_B}$ and
$A_0\mapsto\lambda^{1/4}A_0$.

(c) \emph{rsLoRA.} $\rho_{\mathrm{rs}}(B,A)=\rho_{\mathrm{LoRA}}(\sqrt r\,B,A)$ is an isomorphism of objects. Up to the
hyperparameter torus, rsLoRA can be represented as LoRA at the same $\alpha$ with $\eta_B$ multiplied by $\sqrt r$ under
Adam (by $r$ under SGD), or equally as LoRA with $\alpha$ multiplied by $\sqrt r$ at the same rates. The reading ``same
object, different metric'' is a choice of representative within the orbit. Part (c) concerns a fixed rank; it does not
address rsLoRA's rank-stability prescription, how $\alpha$ should scale with $r$.
\end{restate}

\begin{proof}
Here $\delta(B,A)=BA$ has multidegree $(1,1)$, and $\sigma_B=0$ means $B_0=0$, which every $c_B$ leaves unchanged.

(a) Under Adam, $(c_A,c_B)$ acts by $s\mapsto s/(c_Ac_B)$, $\eta_A\mapsto c_A\eta_A$, $\eta_B\mapsto c_B\eta_B$ and
$\sigma_A\mapsto c_A\sigma_A$. In logarithms this is translation by the column space of
\[a=\begin{pmatrix}-1&-1\\1&0\\0&1\\1&0\end{pmatrix}\qquad(\text{rows } \log s,\log\eta_A,\log\eta_B,\log\sigma_A),\]
which has rank $2$, so the orbit space has dimension $2$ and a basis of the vectors $b$ with $b^\top a=0$ gives a complete
set of invariant monomials (\bgref{D.20}). The vectors $(1,1,1,0)$ and $(0,-1,0,1)$ satisfy $b^\top a=0$ and are independent;
they give $s\eta_A\eta_B$ and $\sigma_A/\eta_A$. Since $\alpha\propto s$ at fixed $r$, Schulman et al.'s quantities are
$s\sigma_A\eta_B=(s\eta_A\eta_B)(\sigma_A/\eta_A)$ and $\sigma_A/\eta_A$, with exponent vectors $(1,0,1,1)$ and $(0,-1,0,1)$,
another basis of the same lattice of invariant monomials. Under SGD the rates scale by $c_A^2$ and $c_B^2$; the weight
matrix has rows $(-1,-1)$, $(2,0)$, $(0,2)$, $(1,0)$, again of rank $2$, and $(2,1,1,0)$ and $(0,-1,0,2)$ give the
invariants $s^2\eta_A\eta_B$ and $\sigma_A^2/\eta_A$.

(b) LoRA+ has $\eta_B=\lambda\eta_A$. Under Adam apply the torus element $c_A=1$, $c_B=1/\lambda$: it maps
$(s,\eta_A,\lambda\eta_A,\sigma_A)$ to $(\lambda s,\eta_A,\eta_A,\sigma_A)$, plain LoRA at scale $\lambda s$ with the same
$\eta_A$ and the same $A_0$; a shared schedule $\varphi(t)$ stays shared. \thmref{Theorem}{III.10} gives
$\varepsilon_B'=\varepsilon/c_B=\lambda\varepsilon$ and $\mathrm{wd}_B'=\mathrm{wd}/c_B=\lambda\cdot\mathrm{wd}$, which are
the stated per-block values, and both vanish when $\varepsilon=\mathrm{wd}=0$. Under SGD, $c_A=1$ and $c_B=\lambda^{-1/2}$
give $\eta_B'=c_B^2\lambda\eta_A=\eta_A$ and $s'=\sqrt\lambda\,s$. The single-rate form uses $c_A=\lambda^{1/4}$ and
$c_B=\lambda^{-1/4}$: then $\eta_A'=\sqrt\lambda\,\eta_A$, $\eta_B'=\lambda^{-1/2}\lambda\eta_A=\sqrt\lambda\,\eta_A$, both
equal to $\sqrt{\eta_A\eta_B}$, $s'=s$, and $\sigma_A'=\lambda^{1/4}\sigma_A$, that is, $A_0\mapsto\lambda^{1/4}A_0$.

(c) $\rho_{\mathrm{rs}}(B,A)=W_0+(\alpha/\sqrt r)BA=W_0+(\alpha/r)(\sqrt rB)A=\rho_{\mathrm{LoRA}}(\sqrt rB,A)$, and
$(B,A)\mapsto(\sqrt rB,A)$ is a pointed diffeomorphism ($0\mapsto0$ in the $B$-factor), hence an isomorphism of objects.
The rsLoRA scale $s_{\mathrm{rs}}=\alpha/\sqrt r$ is the LoRA scale for $\alpha'=\sqrt r\,\alpha$, which is the second
representation. Applying the torus element $c_A=1$, $c_B=\sqrt r$ to rsLoRA gives $s'=s_{\mathrm{rs}}/\sqrt r=\alpha/r$,
the LoRA scale at the same $\alpha$, with $\eta_B'=\sqrt r\,\eta_B$ under Adam and $\eta_B'=r\eta_B$ under SGD, and with
$\varepsilon_B'=\varepsilon/\sqrt r$ and $\mathrm{wd}_B'=\mathrm{wd}/\sqrt r$ when these are nonzero. Both representations
lie on the torus orbit of rsLoRA's hyperparameters, which is what ``up to the hyperparameter torus'' means. The rank $r$
is fixed throughout.
\end{proof}

% ====================================================================================================
\subsection{Theorem III.12}\label{proof:III.12}

\begin{restate}{Theorem}{III.12}{the optimizer gauge hierarchy}
On LoRA's full-rank locus, with hyperparameters shared by all rank channels, the largest subgroup of the gauge $\GL_r$
under which an update rule is equivariant, at the level of parameters, is as follows.

(a) SGD and GD with any block rates, with or without heavy-ball momentum: $\Ort(r)$.

(b) Adam and AdamW at fixed learning rates ($\beta_1,\beta_2\in[0,1)$, any $\varepsilon\ge0$, any weight decay): the
signed permutations $B_r$, of order $2^rr!$. A positive diagonal $g=\diag(d)$ is a symmetry only \emph{jointly} with
per-rank-channel hyperparameters, and only without global-norm clipping. At $\varepsilon=0$ and without weight decay it
suffices to rescale the rates, $\eta_{B,j}\mapsto\eta_{B,j}/d_j$ and $\eta_{A,j}\mapsto d_j\eta_{A,j}$. With
$\varepsilon>0$ one must also give each channel its own $\varepsilon_{B,j}=d_j\varepsilon$ and
$\varepsilon_{A,j}=\varepsilon/d_j$. With decoupled decay multiplied by the learning rate (PyTorch AdamW) one must also
set $\lambda_{B,j}=d_j\lambda$ and $\lambda_{A,j}=\lambda/d_j$; under the convention of \citet{xloshchilov2019decoupled},
where decay is scaled only by the schedule multiplier, $\lambda$ needs no change.

(c) Undamped scaled GD ($\delta=0$; the Riemannian preconditioner of Tong, Ma and Chi and of Zhang and Pilanci, with or
without heavy-ball momentum): all of $\GL_r$. The $W$-trajectory from gauge-related initial states (with zero or
correspondingly transformed momentum) is independent of the representative; a step may drop rank, so the iteration maps
into $W_0+\Mr$.
\end{restate}

\begin{proof}
We use the conventions of \S\ref{proof:III.1}: at $g\cdot q$ the gradients are \eqref{eq:C2-gauge-grad}, and equivariance
under $g$ means \eqref{eq:C2-equivariant-step} for all gradients, together with transported optimizer state. The set of
such $g$ is a subgroup. On the full-rank locus $\nabla_B=sGA^\top$ ranges over all of $\R^{m\times r}$ as $G$ ranges over
$\R^{m\times n}$, because $A$ has full row rank; this is the only use of the locus in the inclusions $\subseteq$ below.

(a) An SGD step is $\delta B=-\eta_B\nabla_B$, $\delta A=-\eta_A\nabla_A$. At $g\cdot q$,
$\delta B'=-\eta_B\nabla_Bg^\top$, and \eqref{eq:C2-equivariant-step} requires $\nabla_B(g^\top-g^{-1})=0$ for all
$\nabla_B\in\R^{m\times r}$, that is, $g^\top=g^{-1}$. Conversely, for $g\in\Ort(r)$, $\delta B'=\delta Bg^\top=\delta Bg^{-1}$
and $\delta A'=-\eta_Ag^{-\top}\nabla_A=g\,\delta A$. A heavy-ball or Nesterov buffer is a linear combination of past
gradients, all taken at points related by the same $g$, so for orthogonal $g$ it transforms like the gradients, $b'_B=b_Bg^{-1}$
and $b'_A=gb_A$, and the step is transported. The first step from a zero buffer is a positive multiple of an SGD step (the multiple is $1$ for heavy ball
and $1+\mu_m$ for Nesterov with momentum coefficient $\mu_m$), so the inclusion $\subseteq\Ort(r)$ also holds with momentum.

(b) \emph{The inclusion $\supseteq$.} Let $g$ be a signed permutation, so $g^{-1}=g^\top$ and $\nabla_{B'}=\nabla_Bg^{-1}$
permutes the columns of $\nabla_B$ and flips some signs; likewise $\nabla_{A'}=g\nabla_A$ permutes and flips rows. Adam is
coordinatewise: the first moment $m$ is permuted with the same sign flips, the second moment $v$ is permuted without sign
changes, so the normalised step is transported, and so is a decoupled or coupled decay term, which is linear in the
parameters. Hence the step satisfies \eqref{eq:C2-equivariant-step}. Since $g$ is orthogonal it preserves
$\lVert\nabla_B\rVert$ and $\lVert\nabla_A\rVert$, so global-norm clipping does not affect this inclusion.

\emph{The inclusion $\subseteq$.} Consider the first step from zero state at a full-rank point. With bias correction it is
$\delta B=-\eta\big(\phi(\nabla_B)+\lambda B\big)$ with $\phi(x)=x/(\lvert x\rvert+\varepsilon)$ entrywise and $0/0=0$;
without bias correction $\phi$ is replaced by a positive multiple of $x/(\lvert x\rvert+\varepsilon')$ with a rescaled
$\varepsilon'$, which changes nothing below. At $g\cdot q$ the decay term $\lambda Bg^{-1}$ is already transported, so
equivariance requires $\phi(\nabla_Bg^\top)=\phi(\nabla_B)g^{-1}$, row by row:
\begin{equation}\label{eq:C2-adam-row}
\phi(xg^\top)=\phi(x)h\qquad\text{for every row vector }x\in\R^{1\times r},\qquad h:=g^{-1}.
\end{equation}
(With coupled decay the argument of $\phi$ is $\nabla_B+\lambda B$; letting $\nabla_B$ grow along a fixed direction leads
to the limit \eqref{eq:C2-adam-sign} below in the same way.)

\emph{Case $\varepsilon=0$.} Then $\phi=\operatorname{sign}$ and \eqref{eq:C2-adam-row} reads
\begin{equation}\label{eq:C2-adam-sign}
\operatorname{sign}(xg^\top)=\operatorname{sign}(x)h\qquad\text{for every }x.
\end{equation}
Fix $\mathsf s\in\{\pm1\}^r$. The open orthant $\{x:\operatorname{sign}(x)=\mathsf s\}$ is not contained in the finite union
of the hyperplanes $\{x:(xg^\top)_j=0\}$, which are proper because $g$ is invertible (a finite union of hyperplanes has
empty interior), so it contains an $x$ with $xg^\top$
free of zeros. For that $x$, \eqref{eq:C2-adam-sign} gives $\mathsf sh\in\{\pm1\}^r$. Thus $\mathsf sh\in\{\pm1\}^r$ for
every $\mathsf s$. Flipping $\mathsf s_i$ changes $(\mathsf sh)_j$ by $2\mathsf s_ih_{ij}$, a difference of two elements of
$\{\pm1\}$, so $h_{ij}\in\{0,\pm1\}$. If column $j$ of $h$ had no nonzero entry, $(\mathsf sh)_j=0$; if it had $t\ge2$
nonzero entries, the choice $\mathsf s_i=h_{ij}$ on its support would give $(\mathsf sh)_j=t\ge2$. So every column of $h$
has exactly one entry, equal to $\pm1$, and these entries lie in distinct rows because $h$ is invertible. Hence $h$, and
with it $g=h^{-1}$, is a signed permutation matrix.

\emph{Case $\varepsilon>0$.} For a fixed $x$ and $\tau\to\infty$, $\phi(\tau y)\to\operatorname{sign}(y)$ entrywise for every
$y$, so applying \eqref{eq:C2-adam-row} to $\tau x$ and letting $\tau\to\infty$ gives \eqref{eq:C2-adam-sign}, and the case
$\varepsilon=0$ shows $g\in B_r$. (For $\tau\to0$, $\phi(\tau y)=\tau y/\varepsilon+O(\tau^2)$, and the same comparison
gives $g^\top=h$, that is, $g\in\Ort(r)$, which is consistent.) The order of $B_r$ is $2^rr!$: $r!$ permutations times
$2^r$ sign patterns.

\emph{The joint torus.} Let $g=\diag(d)$ with $d_j>0$, and give column $j$ of $B$ the hyperparameters
$(\eta'_{B,j},\varepsilon'_{B,j},\lambda'_{B,j})$ and row $j$ of $A$ the hyperparameters $(\eta'_{A,j},\varepsilon'_{A,j},\lambda'_{A,j})$.
By \eqref{eq:C2-gauge-grad}, column $j$ of $\nabla_{B'}$ is $d_j$ times column $j$ of $\nabla_B$, while
\eqref{eq:C2-equivariant-step} asks for column $j$ of $\delta B'$ to be $1/d_j$ times that of $\delta B$. Suppose that before
a step $B'_{:,j}=B_{:,j}/d_j$, $m'_{:,j}=d_jm_{:,j}$ and $v'_{:,j}=d_j^2v_{:,j}$ (true for zero initial state). After the
moment updates these relations persist, and the normalised step of column $j$ is
\[\frac{d_j\hat m}{d_j\sqrt{\hat v}+\varepsilon'_{B,j}}=\frac{\hat m}{\sqrt{\hat v}+\varepsilon'_{B,j}/d_j},\]
which equals the original normalised step iff $\varepsilon'_{B,j}=d_j\varepsilon$ (at $\varepsilon=0$ both are $0$). The
update of column $j$ is then $-\eta'_{B,j}\big(\text{step}+\lambda'_{B,j}B_{:,j}/d_j\big)$, which equals
$-\tfrac1{d_j}\eta\big(\text{step}+\lambda B_{:,j}\big)$ iff $\eta'_{B,j}=\eta/d_j$ and $\eta'_{B,j}\lambda'_{B,j}/d_j=\eta\lambda/d_j$,
that is, $\lambda'_{B,j}=d_j\lambda$. Under the convention of \citet{xloshchilov2019decoupled} the decay term is
$-\zeta_t\lambda'_{B,j}B'_{:,j}$ with the schedule multiplier $\zeta_t$ alone, and it is transported iff $\lambda'_{B,j}=\lambda$.
Row $j$ of $\nabla_{A'}$ is $1/d_j$ times that of $\nabla_A$ and the step must be multiplied by $d_j$, which gives
$\varepsilon'_{A,j}=\varepsilon/d_j$, $\eta'_{A,j}=d_j\eta$ and $\lambda'_{A,j}=\lambda/d_j$ (no change under
Loshchilov--Hutter). The relations persist by induction. With channel-uniform hyperparameters a positive diagonal
$g\ne I$ is not in $B_r$, so by the inclusion $\subseteq$ it is a symmetry only jointly with these per-channel values.
The induction compares the two runs at equal clipping factors. Global-norm clipping multiplies all gradients by
$\min(1,C/\lVert(\nabla_B,\nabla_A)\rVert)$, and
$\lVert(\nabla_{B'},\nabla_{A'})\rVert^2=\sum_jd_j^2\lVert\nabla_{B,:,j}\rVert^2+\sum_jd_j^{-2}\lVert\nabla_{A,j,:}\rVert^2$
differs from the original norm for $d\ne\mathbf 1$ in general; the two runs then weight their gradient histories
differently, and the relations above fail. This is why the joint symmetry is stated without global-norm clipping.

(c) By \thmref{Proposition}{III.4}(a) the undamped scaled-GD step satisfies \eqref{eq:C2-equivariant-step} for every
$g\in\GL_r$. With heavy-ball momentum on the preconditioned gradients, $p_B=\nabla_B(AA^\top)^{-1}$ transforms as
$p_{B'}=p_Bg^{-1}$, so a buffer of such terms transforms in the same way and so does the step. With momentum on the raw
gradients followed by preconditioning, the buffer transforms as $b_{B'}=b_Bg^\top$ because all past points are related by
the same $g$, and $b_{B'}(A'A'^\top)^{-1}=b_Bg^\top g^{-\top}(AA^\top)^{-1}g^{-1}$ is again transported. The $A$-factor is
symmetric. So from gauge-related initial states, with zero or correspondingly transformed momentum, $q'_t=g\cdot q_t$ for
all $t$ and the $W$-trajectories coincide. Each iterate lies in $W_0+\Mr$, and the iteration is defined as long as both
factors keep full rank; a step can lower the rank of $B+\delta B$ or $A+\delta A$, so the map goes into $W_0+\Mr$ rather
than $W_0+\MM_r$.

The hierarchy $B_r\subset\Ort(r)\subset\GL_r$ is strict for $r\ge2$, since $B_r$ is finite, $\Ort(r)$ is infinite and
compact, and $\GL_r$ is not compact.
\end{proof}

\begin{proof}[Proof of the statements about published variants]
\emph{Damped scaled GD.} With a damping constant $\delta>0$ the two inverses become $(AA^\top+\delta I)^{-1}$ and
$(B^\top B+\delta I)^{-1}$.
Equivariance of the $B$-step requires, for all $\nabla_B$,
\[\nabla_Bg^\top(gAA^\top g^\top+\delta I)^{-1}=\nabla_B(AA^\top+\delta I)^{-1}g^{-1},\quad\text{that is,}\quad
g^\top(gAA^\top g^\top+\delta I)^{-1}=(AA^\top+\delta I)^{-1}g^{-1}.\]
Inverting both sides gives $(gAA^\top g^\top+\delta I)g^{-\top}=g(AA^\top+\delta I)$, that is,
$gAA^\top+\delta g^{-\top}=gAA^\top+\delta g$, so $g^{-\top}=g$ and $g\in\Ort(r)$. Conversely, for orthogonal $g$ we have
$(gAA^\top g^\top+\delta I)^{-1}=g(AA^\top+\delta I)^{-1}g^\top$, and both steps are transported. The computation does not use full rank, so it holds at every point, including $B=0$.

\emph{Scaled AdamW} applies Adam to the preconditioned gradient $p_B=\nabla_B(AA^\top+\delta I)^{-1}$ (and similarly for
$A$), with $\delta\ge0$. For a signed permutation $g$, $p_{B'}=p_Bg^{-1}$ by the previous paragraph, and Adam commutes with
signed permutations, so the rule is $B_r$-equivariant. Conversely, the first step requires
$\phi(\nabla_BM')=\phi(\nabla_BM)g^{-1}$ for all $\nabla_B$, with $M=(AA^\top+\delta I)^{-1}$ and
$M'=g^\top(gAA^\top g^\top+\delta I)^{-1}$ invertible. Writing $y=xM$ and letting $y$ grow, as in the case $\varepsilon>0$
above, gives $\operatorname{sign}(yM^{-1}M')=\operatorname{sign}(y)g^{-1}$ for all $y$, and the orthant argument, which
used only invertibility of the matrix on the left, gives $g^{-1}\in B_r$.

\emph{LoRA-Pro with the Sylvester-optimal $X$.} Write $\mathcal P=\R^{m\times r}\times\R^{r\times n}$ with its Frobenius
inner product, $u=(u_B,u_A)\in\mathcal P$ for a parameter velocity, and $\mathcal K_q=\ker D\rho_q$. On the full-rank locus
$\mathcal K_q=\{(-BX,XA):X\in\R^{r\times r}\}$: if $\dot BA=-B\dot A$, then $\dot B=-BX$ with $X=\dot AA^\top(AA^\top)^{-1}$,
and $BXA=B\dot A$ gives $\dot A=XA$. By \S\ref{proof:III.4}(b), LoRA-Pro's directions $(g^B,g^A)$ form the affine space
$\mathcal M_q(G)=\{u:D\rho_qu=\Pi_{\mathcal T}G\}$, and LoRA-Pro chooses $X$ by minimising the Frobenius distance from
$(g^B,g^A)$ to LoRA's own gradients $D\rho_q^\top G=(sGA^\top,sB^\top G)$; the minimiser solves a Sylvester equation
(\bgref{D.33}, \auxref{B.20}). Since $D\rho_q^\top G\perp\mathcal K_q$, for $u\in\mathcal M_q(G)$ the inner product
$\langle u,D\rho_q^\top G\rangle$ is constant, and minimising $\lVert u-D\rho_q^\top G\rVert^2$ over $\mathcal M_q(G)$ is the
same as minimising $\lVert u\rVert^2$. So LoRA-Pro's direction is the minimum-norm solution $u_q(G)$ of
$D\rho_qu=\Pi_{\mathcal T}G$, that is, $u_q(G)=D\rho_q^+G$.

We show that $u_{g\cdot q}(G)=\varphi_g(u_q(G))$ for all $G$ iff $g\in\Ort(r)$, where $\varphi_g(u)=(u_Bg^{-1},gu_A)$ is
the linear action. Since $D\rho_{g\cdot q}\varphi_g=D\rho_q$ and $\Pi_{\mathcal T}$ depends only on $W$, substituting
$u=\varphi_g(w)$ shows that $u_{g\cdot q}(G)=\varphi_g(w^*)$, where $w^*$ minimises $\langle w,\mathcal Nw\rangle$ over
$\mathcal M_q(G)$ with $\mathcal N=\varphi_g^\top\varphi_g$. The adjoint of $\varphi_g$ is $\varphi_g^\top(v_B,v_A)=(v_Bg^{-\top},g^\top v_A)$,
since $\langle w_Bg^{-1},v_B\rangle=\langle w_B,v_Bg^{-\top}\rangle$ and $\langle gw_A,v_A\rangle=\langle w_A,g^\top v_A\rangle$, so
$\mathcal N(w_B,w_A)=(w_Bg^{-1}g^{-\top},g^\top gw_A)=(w_B(g^\top g)^{-1},(g^\top g)w_A)$. So
equivariance holds iff, on every coset of $\mathcal K_q$ (as $G$ varies, $\mathcal M_q(G)$ runs through all cosets, since
$D\rho_q$ maps onto $\mathcal T$), the $\mathcal N$-minimum-norm point equals the Euclidean one. The two minimisers are the
intersections of the coset with $\mathcal K_q^{\perp_{\mathcal N}}=\mathcal N^{-1}(\mathcal K_q^\perp)$ and with
$\mathcal K_q^\perp$; they agree on every coset iff these two complements are equal (each point of either complement is the minimiser on its
own coset), iff $\mathcal N\mathcal K_q=\mathcal K_q$
($\mathcal N$ is symmetric and invertible). Put $H=g^\top g\succ0$. Then $\mathcal N(-BX,XA)=(-BXH^{-1},HXA)$ lies in
$\mathcal K_q$ iff some $Y$ has $BY=BXH^{-1}$ and $YA=HXA$, that is (as $B$ is injective and $A$ has full row rank),
$XH^{-1}=HX$. Requiring this for every $X\in\R^{r\times r}$, the matrix-unit lemma of \S\ref{proof:III.3} gives
$H^{-1}=cI=H$, so $c^2=1$, and $c=1$ because $H$ is positive definite. Hence $g^\top g=I$ and $g\in\Ort(r)$.
Conversely, orthogonal $g$ give $\mathcal N=I$. So the parameter update with the Sylvester-optimal $X$ is equivariant
exactly under $\Ort(r)$.

The variant with $X=0$ is $g^A=s^{-1}(B^\top B)^{-1}B^\top G$, $g^B=s^{-1}(I-P_U)GA^\top(AA^\top)^{-1}$. At $g\cdot q$,
$P_U$ is unchanged and
\[g^{A'}=s^{-1}g(B^\top B)^{-1}g^\top g^{-\top}B^\top G=g\,g^A,\qquad g^{B'}=s^{-1}(I-P_U)GA^\top g^\top g^{-\top}(AA^\top)^{-1}g^{-1}=g^Bg^{-1},\]
so it is $\GL_r$-equivariant and its discrete $W$-step descends exactly. The first-order $W$-step of either variant is
$-\eta\Pi_{\mathcal T}G$ (\thmref{Proposition}{III.4}(b)). For the Sylvester variant, $\varphi_g^{-1}u_{g\cdot q}(G)-u_q(G)\in\mathcal K_q$
equals $(-BY,YA)$ for some $Y$, and the second-order term $s\,\delta B\,\delta A=\eta^2s\,g^Bg^A$ becomes
$\eta^2s(g^B-BY)(g^A+YA)$ at $g\cdot q$. The difference $\eta^2s(g^BYA-BYg^A-BY^2A)$ is exactly $\eta^2$ times a matrix that does not depend on $\eta$, and it
is nonzero in general. Numerically, at random
full-rank points with $(m,n,r)=(7,6,3)$, $s=0.8$ and $g$ a Gaussian matrix plus $2I$, the largest entry of the
difference of the two $W$-steps at $\eta=0.1$ ranges from $8\times10^{-4}$ to $1.2\times10^{-2}$ over seven random draws,
and it falls by a factor of exactly $100$ each time $\eta$ falls by $10$. It is of
order $10^{-15}$ (rounding) for an orthogonal $g$ and for the variant with $X=0$.

\emph{LoRA-RITE.} In the notation of \citet{yen2024lora} the update is $AB^\top$ with $A\in\R^{m\times r}$ and
$B\in\R^{n\times r}$, and the gauge is $A_2=A_1R$, $B_2=B_1R^{-\top}$ with $R\in\GL_r$. We follow their Algorithm~1 line
by line on the full-rank locus, where its pseudo-inverses are inverses. The rule is symmetric in $A$ and $B$, so we treat
the step of $A$. Algorithm~1 writes $B=U_BR_B$ by polar decomposition, with $U_B$ having orthonormal columns, and
forms, at step $t$:
\begin{itemize}
\item the unmagnified gradient $\bar\nabla_{A,t}=\nabla_{A,t}R_{B_t}^{-1}$ and the basis change
$P_{A,t}=U_{B_t}^\top U_{B_{t-1}}$;
\item a second moment $\bar V_{A,t}$, a combination with fixed weights of $P_{A,t}\bar V_{A,t-1}P_{A,t}^\top$ and
$\bar\nabla_{A,t}^\top\bar\nabla_{A,t}$;
\item the escaped mass
\[\rho_{A,t}=\rho_{A,t-1}+d_\lambda\big(\bar V_{A,t-1},P_{A,t}\bar V_{A,t-1}P_{A,t}^\top\big),\qquad
d_\lambda(E_1,E_2)=\max_i\lvert\lambda_i(E_1)-\lambda_i(E_2)\rvert,\]
which depends only on the spectra of its two arguments;
\item the preconditioned step, the first moment and the step
\[\bar S_{A,t}=\bar\nabla_{A,t}(\bar V_{A,t}+\rho_{A,t}I)^{-1/2},\quad
\bar M_{A,t}=\beta_1\bar M_{A,t-1}P_{A,t}^\top+(1-\beta_1)\bar S_{A,t},\quad
\delta A_t=-\eta\,\bar M_{A,t}R_{B_t}^{-\top}.\]
\end{itemize}
Run the algorithm from $(A_1,B_1)$ and from $(A_2,B_2)=(A_1R,B_1R^{-\top})$, with zero initial moments and escaped
mass. Suppose that at step $t$ the parameters of the two runs are related by $R$. Then $B_{2,t}$ and $B_{1,t}$ have the
same column space, so $U_{B_{2,t}}=U_{B_{1,t}}O_t$ for some $O_t\in\Ort(r)$, and
$R_{B_{2,t}}=U_{B_{2,t}}^\top B_{2,t}=O_t^\top R_{B_{1,t}}R^{-\top}$. With $G$ the gradient in the weight,
$\nabla_{A_2}=GB_2=\nabla_{A_1}R^{-\top}$, so
\[\bar\nabla_{A_2,t}=\nabla_{A_1,t}R^{-\top}R^\top R_{B_{1,t}}^{-1}O_t=\bar\nabla_{A_1,t}O_t.\]
Suppose also that the internal states are related by $\bar V^{(2)}_{t-1}=O_{t-1}^\top\bar V^{(1)}_{t-1}O_{t-1}$,
$\bar M^{(2)}_{t-1}=\bar M^{(1)}_{t-1}O_{t-1}$ and $\rho^{(2)}_{t-1}=\rho^{(1)}_{t-1}$; at $t=1$ this holds because all
three start at zero. Then $P^{(2)}_{A,t}=O_t^\top P^{(1)}_{A,t}O_{t-1}$, so
$P^{(2)}_{A,t}\bar V^{(2)}_{t-1}P^{(2)\top}_{A,t}=O_t^\top P^{(1)}_{A,t}\bar V^{(1)}_{t-1}P^{(1)\top}_{A,t}O_t$ and
$\bar V^{(2)}_t=O_t^\top\bar V^{(1)}_tO_t$. Conjugation by an orthogonal matrix keeps eigenvalues, so
$\rho^{(2)}_t=\rho^{(1)}_t$, and it commutes with $(\,\cdot+\rho I)^{-1/2}$, so $\bar S^{(2)}_t=\bar S^{(1)}_tO_t$. Since
$\bar M^{(2)}_{t-1}P^{(2)\top}_{A,t}=\bar M^{(1)}_{t-1}P^{(1)\top}_{A,t}O_t$, also $\bar M^{(2)}_t=\bar M^{(1)}_tO_t$.
Finally $R_{B_{2,t}}^{-\top}=O_t^\top R_{B_{1,t}}^{-\top}R$, so
\[\delta A_{2,t}=-\eta\,\bar M^{(1)}_tO_tO_t^\top R_{B_{1,t}}^{-\top}R=\delta A_{1,t}R,\]
the transported step. The step of $B$ is the same computation with $A$ and $B$ exchanged and $R$ replaced by
$R^{-\top}$, so $\delta B_{2,t}=\delta B_{1,t}R^{-\top}$. The new parameters are again related by $R$, and the new
internal states by the orthogonal matrices $O_t$, which completes the induction. The resulting step agrees with Eq.~15
of \citet{yen2024lora}, $\delta A_2=\delta A_1B_1^\top(B_2^\top)^+$, since $(B_2^\top)^+=B_1(B_1^\top B_1)^{-1}R$ when
$B_1$ has full column rank, and with the numerical replication reported after \thmref{Theorem}{III.12}, which agrees
to $3\times10^{-14}$. LoRA-RITE is
therefore an Adam-type optimizer in class (c).
\end{proof}

% ====================================================================================================
\subsection{Corollary III.13}\label{proof:III.13}

\begin{restate}{Corollary}{III.13}{what a zero-$B$ initialisation really chooses}
At $q_0=(0,A_0)$ with $\rk A_0=r$, with hyperparameters shared by the rank channels and the same data order and dropout
masks, two choices of $A_0$ give the same $\theta$-trajectory whenever the optimizer's gauge group relates them. The
resulting invariant of $A_0$ is:
\begin{itemize}
\item $\operatorname{row}A_0\in\Gr(r,n)$ for an update rule that is $\GL_r$-equivariant at every point of its trajectory,
including $B=0$ and rank-deficient $B$, run without global-norm clipping. An example is
$\delta B=-\eta\nabla_B(AA^\top)^{-1}$, $\delta A=-\eta(AA^\top)\nabla_A$, which needs only $\rk A=r$;
\item the PSD matrix $A_0^\top A_0$ for SGD, which is exactly the initial preconditioner in
$\dot W(0)=-s^2\eta_BGA_0^\top A_0$;
\item $A_0$ up to signed permutation of its rows for Adam.
\end{itemize}
The undamped scaled-GD optimizers of \thmref{Theorem}{III.12}(c) are not defined at $B=0$, since they need
$(B^\top B)^{-1}$. LoRA-RITE uses pseudo-inverses and so is defined there, but its $\GL_r$-equivariance is established only
on the full-rank locus. Their damped forms are only $\Ort(r)$-equivariant, so in practice their invariant is a
$\delta$-dependent function of $A_0^\top A_0$, and their Adam-wrapped forms (Riemannian-preconditioned AdamW, HF's LoRA-FA
optimizer) have the invariant $A_0$ modulo $B_r$. Conversely, different invariants give different dynamics, for SGD and
for Adam at $\varepsilon=0$.

Random-$A$ LoRA and EVA are different points of the $B=0$ stratum $\{(0,A):\rk A=r\}$ of the apex fibre
$\rho^{-1}(\theta_0)=\{(B,A):BA=0\}$; Init[B] lies in another stratum of the same fibre. LoRA-FA is a different
\emph{object}, the sub-object $\mathrm{FA}(A_0)\hookrightarrow\mathrm{LoRA}_{(0,A_0)}$ (\thmref{Proposition}{I.8}). More
precisely, EVA is a point of the apex fibre of $\mathrm{LoRA}_{r_l}$ at its per-layer redistributed rank and adjusted
$\alpha$, with orthonormal rows only when \texttt{whiten=False}.
\end{restate}

\begin{proof}
\emph{Equivariance at $B=0$.} Only the inclusions $\supseteq$ of \thmref{Theorem}{III.12} are used. For SGD and $g\in\Ort(r)$
the identities $\delta B'=-\eta_B\nabla_Bg^\top=\delta Bg^{-1}$ and $\delta A'=g\,\delta A$ hold at every point, with no rank
condition, and similarly for momentum buffers. For Adam and $g\in B_r$ the coordinatewise argument also holds at every
point, and the initial state $(m,v)=(0,0)$ is fixed by $B_r$. Orthogonal $g$ preserve $\lVert\nabla_B\rVert$ and
$\lVert\nabla_A\rVert$ by \eqref{eq:C2-gauge-grad}, so global-norm clipping factors agree in the two runs. The start
transforms as $g\cdot(0,A_0)=(0,gA_0)$, since $0\cdot g^{-1}=0$. By the induction of \S\ref{proof:III.1}, $A_0$ and
$gA_0$ give the same $\theta$-trajectory whenever $g$ lies in the optimizer's group. In the first bullet the hypothesis is
equivariance under $\GL_r$ at every point of the trajectory, and the same induction applies.

\emph{The example rule.} At $g\cdot q$, by \eqref{eq:C2-gauge-grad} and
$(A'A'^\top)^{-1}=g^{-\top}(AA^\top)^{-1}g^{-1}$,
\[\delta B'=-\eta\nabla_Bg^\top g^{-\top}(AA^\top)^{-1}g^{-1}=\delta B\,g^{-1},\qquad
\delta A'=-\eta\,gAA^\top g^\top g^{-\top}\nabla_A=g\,\delta A .\]
Only $\rk A=r$ was used, so $B$ may be $0$ or rank deficient. Since $\delta A=-\eta AA^\top sB^\top G=A\,N$ with $N=-\eta sA^\top B^\top G$, the step maps $A$ to $A(I+N)$, which
keeps rank $r$ whenever $I+N$ is invertible, in particular for small $\eta$; in continuous time $\dot A=AN(t)$ gives
$A(t)=A_0\Psi(t)$ with $\Psi$ an invertible fundamental matrix, so $\rk A(t)=r$ for all $t$.

\emph{The invariants.} The invariant of $A_0$ is its orbit under the optimizer's group acting by $A_0\mapsto gA_0$. For
$\Ort(r)$, $A_0'=OA_0$ for some $O\in\Ort(r)$ iff $A_0'^\top A_0'=A_0^\top A_0$ (\thmref{Theorem}{II.5}(2)). For $B_r$ the
orbit is the set of signed row permutations of $A_0$. For $\GL_r$, two full-row-rank matrices satisfy $A_0'=gA_0$ for some
(unique) $g$ iff they have the same row space (proof of \thmref{Theorem}{II.2}), so the invariant is
$\operatorname{row}A_0\in\Gr(r,n)$ (\bgref{C.15}). At $B_0=0$ the LoRA cometric of \thmref{Observation}{III.2} is
$K_{(0,A_0)}(G)=s^2\eta_BGA_0^\top A_0$, so $\dot W(0)=-s^2\eta_BGA_0^\top A_0$.

\emph{Practical forms.} Undamped scaled GD needs $(B^\top B)^{-1}$, which does not exist at $B=0$. The damped forms are
$\Ort(r)$-equivariant at every point (\S\ref{proof:III.12}, published variants), so their trajectories from $(0,A_0)$
depend on $A_0$ only through $A_0^\top A_0$, in a way that depends on $\delta$. A preconditioned step wrapped in Adam(W) is
$B_r$-equivariant at every point (same place), so the invariant is $A_0$ modulo $B_r$.

\emph{Converse for SGD.} The first SGD step from $(0,A_0)$ is $\delta B=-\eta_BsGA_0^\top$, $\delta A=-\eta_AsB_0^\top G=0$,
so $\delta W_1=s\,\delta B\,A_0=-s^2\eta_BGA_0^\top A_0$ exactly. If $A_0$ and $A_0'$ give the same trajectory for every loss,
the linear losses $\langle G,\cdot\rangle$ give $GA_0^\top A_0=GA_0'^\top A_0'$ for every $G$, so $A_0^\top A_0=A_0'^\top A_0'$.
Hence different invariants give different first steps for some loss.

\emph{Converse for Adam at $\varepsilon=0$.} The first step from $(0,A_0)$ with zero state has $\nabla_A=0$, so $\delta A=0$,
and $\delta B=-\eta_B\operatorname{sign}(sGA_0^\top)$, so $\delta W_1=-s\eta_B\operatorname{sign}(GA_0^\top)A_0$ (without bias
correction, a positive multiple of this). If $A_0$ and $A_0'$ give the same trajectory for every loss, then row by row
\[F(x):=\operatorname{sign}(xA_0^\top)A_0=\sum_{j=1}^r\operatorname{sign}(x\cdot a_j)\,a_j\]
equals $F'(x)=\sum_j\operatorname{sign}(x\cdot a'_j)\,a'_j$ for every $x\in\R^{1\times n}$, where $a_j$ and $a'_j$ are the rows
of $A_0$ and $A_0'$. The rows $a_j$ are linearly independent, so the hyperplanes $H_j=a_j^\perp$ are pairwise distinct, and
each $H_j$ contains points $x$ lying on no other $H_k$ (a hyperplane is not covered by finitely many proper subspaces of
it, \auxref{B.36}(b)). At such a point, crossing $H_j$ in the direction of increasing $x\cdot a_j$ changes $F$ by $2a_j\ne0$, while $F$ is
constant on the connected components of the complement of $\bigcup_jH_j$. So the closure of the set of discontinuities of
$F$ is $\bigcup_jH_j$, and likewise for $F'$. From $F=F'$ we get $\bigcup_jH_j=\bigcup_kH'_k$, and since a hyperplane
contained in a finite union of hyperplanes equals one of them (\auxref{B.36}(b) again), there is a permutation $\pi$ with $H_j=H'_{\pi(j)}$, that is,
$a'_{\pi(j)}=c_ja_j$ with $c_j\ne0$. Crossing $H_j$ as above changes $F'$ by
$2\operatorname{sign}(c_j)a'_{\pi(j)}=2\lvert c_j\rvert a_j$, and comparing with $2a_j$ gives $\lvert c_j\rvert=1$. So
$A_0'$ is a signed row permutation of $A_0$. For the example $\GL_r$-rule the same argument is immediate: its first step is
$\delta W_1=-\eta s^2GP_{\operatorname{row}A_0}$, which determines $\operatorname{row}A_0$.

\emph{Strata and the sub-object.} Since $s>0$, $\rho^{-1}(\theta_0)=\{(B,A):W_0+sBA=W_0\}=\{BA=0\}$. Random-$A$ LoRA and
EVA start at $(0,A_0)$ with $\rk A_0=r$, in the stratum $\{(0,A):\rk A=r\}$, at points that differ for a generic random
draw; Init[B] starts at $(B_0,0)$ with $\rk B_0=r$, in the disjoint stratum $\{(B,0):\rk B=r\}$. The map $B\mapsto(B,A_0)$
is injective, pointed ($0\mapsto(0,A_0)$) and commutes with the maps to $\Theta$, since $\rho(B,A_0)=W_0+sBA_0$; after the
rescaling $B\mapsto sB$ its source is $\mathrm{FA}(A_0)$ of \thmref{Proposition}{I.8}. An injective morphism is a
monomorphism, so $\mathrm{FA}(A_0)$ is a sub-object of $\mathrm{LoRA}_{(0,A_0)}$ (\bgref{A.15}). Its image
$W_0+\{XA_0\}$ has dimension $mr$, smaller than the dimension $r(m+n-r)$ of the image $W_0+\Mr$ of
$\mathrm{LoRA}_{(0,A_0)}$ when $r<n$, and images are isomorphism invariants, so the two objects are not isomorphic. The
last sentence describes EVA's construction and needs no proof.
\end{proof}
